\begin{filecontents*}[overwrite]{references.bib}
@article{berlinet2008functional,
  author  = {Berlinet, Alain and Biau, G{\'e}rard and Rouvi{\`e}re, Laurent},
  title   = {Functional Supervised Classification with Wavelets},
  journal = {Annales de l'ISUP},
  volume  = {52},
  number  = {1--2},
  pages   = {61--80},
  year    = {2008}
}

@book{devroye1996probabilistic,
  author    = {Devroye, Luc and Gy{\"o}rfi, L{\'a}szl{\'o} and Lugosi, G{\'a}bor},
  title     = {A Probabilistic Theory of Pattern Recognition},
  publisher = {Springer},
  address   = {New York},
  year      = {1996}
}

@article{devroye1988automatic,
  author  = {Devroye, Luc},
  title   = {Automatic Pattern Recognition: A Study of the Probability of Error},
  journal = {IEEE Transactions on Pattern Analysis and Machine Intelligence},
  volume  = {10},
  number  = {4},
  pages   = {530--543},
  year    = {1988}
}

@article{vapnik1971uniform,
  author  = {Vapnik, Vladimir N. and Chervonenkis, Alexey Ya.},
  title   = {On the Uniform Convergence of Relative Frequencies of Events to Their Probabilities},
  journal = {Theory of Probability and its Applications},
  volume  = {16},
  number  = {2},
  pages   = {264--280},
  year    = {1971}
}

@book{shumway2011time,
  author    = {Shumway, Robert H. and Stoffer, David S.},
  title     = {Time Series Analysis and Its Applications},
  edition   = {3},
  series    = {Springer Texts in Statistics},
  publisher = {Springer},
  address   = {New York},
  year      = {2011}
}

@book{daubechies1992ten,
  author    = {Daubechies, Ingrid},
  title     = {Ten Lectures on Wavelets},
  publisher = {SIAM},
  address   = {Philadelphia},
  year      = {1992}
}

@article{mallat1989theory,
  author  = {Mallat, St{\'e}phane},
  title   = {A Theory for Multiresolution Signal Decomposition: The Wavelet Representation},
  journal = {IEEE Transactions on Pattern Analysis and Machine Intelligence},
  volume  = {11},
  number  = {7},
  pages   = {674--693},
  year    = {1989}
}

@article{haar1910theorie,
  author  = {Haar, Alfr{\'e}d},
  title   = {Zur Theorie der orthogonalen Funktionensysteme},
  journal = {Mathematische Annalen},
  volume  = {69},
  pages   = {331--371},
  year    = {1910}
}

@article{donoho1994ideal,
  author  = {Donoho, David L. and Johnstone, Iain M.},
  title   = {Ideal Spatial Adaptation by Wavelet Shrinkage},
  journal = {Biometrika},
  volume  = {81},
  number  = {3},
  pages   = {425--455},
  year    = {1994}
}

@article{donoho1995wavelet,
  author  = {Donoho, David L. and Johnstone, Iain M. and Kerkyacharian, G{\'e}rard and Picard, Dominique},
  title   = {Wavelet Shrinkage: Asymptopia?},
  journal = {Journal of the Royal Statistical Society, Series B},
  volume  = {57},
  number  = {2},
  pages   = {301--369},
  year    = {1995}
}

@article{hastie1995penalized,
  author  = {Hastie, Trevor and Buja, Andreas and Tibshirani, Robert},
  title   = {Penalized Discriminant Analysis},
  journal = {The Annals of Statistics},
  volume  = {23},
  number  = {1},
  pages   = {73--102},
  year    = {1995}
}

@article{stone1977consistent,
  author  = {Stone, Charles J.},
  title   = {Consistent Nonparametric Regression},
  journal = {The Annals of Statistics},
  volume  = {5},
  number  = {4},
  pages   = {595--645},
  year    = {1977}
}

@article{abraham2006kernel,
  author  = {Abraham, Christophe and Biau, G{\'e}rard and Cadre, Beno{\^\i}t},
  title   = {On the Kernel Rule for Function Classification},
  journal = {Annals of the Institute of Statistical Mathematics},
  volume  = {58},
  pages   = {619--633},
  year    = {2006}
}

@article{cerou2006nearest,
  author  = {C{\'e}rou, Fr{\'e}d{\'e}ric and Guyader, Arnaud},
  title   = {Nearest Neighbor Classification in Infinite Dimension},
  journal = {ESAIM: Probability and Statistics},
  volume  = {10},
  pages   = {340--355},
  year    = {2006}
}

@article{biau2005functional,
  author  = {Biau, G{\'e}rard and Bunea, Florentina and Wegkamp, Marten H.},
  title   = {Functional Classification in {H}ilbert Spaces},
  journal = {IEEE Transactions on Information Theory},
  volume  = {51},
  number  = {6},
  pages   = {2163--2172},
  year    = {2005}
}

@book{ramsay2005functional,
  author    = {Ramsay, James O. and Silverman, Bernard W.},
  title     = {Functional Data Analysis},
  edition   = {2},
  publisher = {Springer},
  address   = {New York},
  year      = {2005}
}

@book{ferraty2006nonparametric,
  author    = {Ferraty, Fr{\'e}d{\'e}ric and Vieu, Philippe},
  title     = {Nonparametric Functional Data Analysis: Theory and Practice},
  publisher = {Springer},
  address   = {New York},
  year      = {2006}
}

@article{wahba1980automatic,
  author  = {Wahba, Grace},
  title   = {Automatic Smoothing of the Log Periodogram},
  journal = {Journal of the American Statistical Association},
  volume  = {75},
  number  = {369},
  pages   = {122--132},
  year    = {1980}
}

@article{moulin1994wavelet,
  author  = {Moulin, Pierre},
  title   = {Wavelet Thresholding Techniques for Power Spectrum Estimation},
  journal = {IEEE Transactions on Signal Processing},
  volume  = {42},
  number  = {11},
  pages   = {3126--3136},
  year    = {1994}
}

@article{hampel1974influence,
  author  = {Hampel, Frank R.},
  title   = {The Influence Curve and Its Role in Robust Estimation},
  journal = {Journal of the American Statistical Association},
  volume  = {69},
  number  = {346},
  pages   = {383--393},
  year    = {1974}
}

@phdthesis{olszewski2001generalized,
  author = {Olszewski, Robert T.},
  title  = {Generalized Feature Extraction for Structural Pattern Recognition in Time-Series Data},
  school = {Carnegie Mellon University},
  year   = {2001}
}

@article{dau2019ucr,
  author  = {Dau, Hoang Anh and Bagnall, Anthony and Kamgar, Kaveh and Yeh, Chin-Chia Michael and Zhu, Yan and Gharghabi, Shaghayegh and Ratanamahatana, Chotirat Ann and Keogh, Eamonn},
  title   = {The {UCR} Time Series Archive},
  journal = {IEEE/CAA Journal of Automatica Sinica},
  volume  = {6},
  number  = {6},
  pages   = {1293--1305},
  year    = {2019}
}

@article{lee2019pywavelets,
  author  = {Lee, Gregory R. and Gommers, Ralf and Waselewski, Filip and Wohlfahrt, Kai and O'Leary, Aaron},
  title   = {{PyWavelets}: A {P}ython Package for Wavelet Analysis},
  journal = {Journal of Open Source Software},
  volume  = {4},
  number  = {36},
  pages   = {1237},
  year    = {2019}
}

@article{pedregosa2011scikit,
  author  = {Pedregosa, Fabian and others},
  title   = {Scikit-learn: Machine Learning in {P}ython},
  journal = {Journal of Machine Learning Research},
  volume  = {12},
  pages   = {2825--2830},
  year    = {2011}
}
\end{filecontents*}

\documentclass[11pt,a4paper,twoside,openright]{report}

% ── Langue et encodage ───────────────────────────────────────────────────────
\usepackage{iftex}
\ifPDFTeX
  \usepackage[T1]{fontenc}
  \usepackage[utf8]{inputenc}
  \usepackage{lmodern}
\else
  \usepackage{fontspec}
\fi
\usepackage[english,french]{babel}
\usepackage{microtype}

% ── Mise en page ─────────────────────────────────────────────────────────────
\usepackage[a4paper,margin=2.6cm,bindingoffset=0.6cm]{geometry}
\usepackage{setspace}
\setstretch{1.12}
\usepackage{fancyhdr}
\pagestyle{fancy}
\fancyhf{}
\fancyhead[LE,RO]{\thepage}
\fancyhead[RE]{\small\nouppercase{\leftmark}}
\fancyhead[LO]{\small\nouppercase{\rightmark}}
\renewcommand{\headrulewidth}{0.3pt}
\setlength{\headheight}{14pt}
\fancypagestyle{plain}{\fancyhf{}\fancyfoot[C]{\thepage}\renewcommand{\headrulewidth}{0pt}}

% ── Mathématiques ────────────────────────────────────────────────────────────
\usepackage{amsmath,amssymb,amsthm,mathtools}
\usepackage{bm}
\usepackage{dsfont}

% ── Tableaux, figures ────────────────────────────────────────────────────────
\usepackage{graphicx}
\graphicspath{{figures/}}
\usepackage{booktabs}
\usepackage{multirow}
\usepackage{array}
\usepackage{caption}
\captionsetup{font=small,labelfont=bf,format=hang}
\usepackage{subcaption}
\usepackage{float}

% ── Divers ───────────────────────────────────────────────────────────────────
\usepackage{enumitem}
\setlist{itemsep=2pt,topsep=4pt}
\usepackage{xcolor}
\definecolor{bleuamu}{HTML}{1F3A68}
\usepackage[natbib=true,style=authoryear,maxcitenames=2,maxbibnames=10,backend=bibtex,uniquename=false,uniquelist=false]{biblatex}
\addbibresource{references.bib}
\DefineBibliographyExtras{french}{\restorecommand\mkbibnamefamily}
\usepackage{csquotes}
\usepackage{xurl}
\usepackage[colorlinks=true,linkcolor=bleuamu,citecolor=bleuamu,urlcolor=bleuamu,hypertexnames=false]{hyperref}
\usepackage[capitalise,noabbrev,nameinlink]{cleveref}
\crefname{theoreme}{Théorème}{Théorèmes}
\crefname{proposition}{Proposition}{Propositions}
\crefname{lemme}{Lemme}{Lemmes}
\crefname{corollaire}{Corollaire}{Corollaires}
\crefname{definition}{Définition}{Définitions}
\crefname{hypothese}{Hypothèse}{Hypothèses}
\crefname{remarque}{Remarque}{Remarques}
\crefname{exemple}{Exemple}{Exemples}
\crefname{equation}{}{}
\crefname{figure}{Figure}{Figures}
\crefname{table}{Tableau}{Tableaux}
\crefname{chapter}{Chapitre}{Chapitres}
\crefname{section}{Section}{Sections}
\crefname{appendix}{Annexe}{Annexes}

% ── Environnements ───────────────────────────────────────────────────────────
\theoremstyle{plain}
\newtheorem{theoreme}{Théorème}[chapter]
\newtheorem{proposition}[theoreme]{Proposition}
\newtheorem{lemme}[theoreme]{Lemme}
\newtheorem{corollaire}[theoreme]{Corollaire}
\theoremstyle{definition}
\newtheorem{definition}[theoreme]{Définition}
\newtheorem{hypothese}{Hypothèse}
\renewcommand{\thehypothese}{(H\arabic{hypothese})}
\newtheorem{exemple}[theoreme]{Exemple}
\theoremstyle{remark}
\newtheorem{remarque}[theoreme]{Remarque}
\renewcommand{\qedsymbol}{$\blacksquare$}

% ── Notations ────────────────────────────────────────────────────────────────
\newcommand{\R}{\mathbb{R}}
\newcommand{\N}{\mathbb{N}}
\newcommand{\PP}{\mathbb{P}}
\newcommand{\E}{\mathbb{E}}
\newcommand{\1}{\mathds{1}}
\newcommand{\Lp}{L^2([0,1])}
\newcommand{\cH}{\mathcal{H}}
\newcommand{\cF}{\mathcal{F}}
\newcommand{\cD}{\mathcal{D}}
\newcommand{\cC}{\mathcal{C}}
\newcommand{\cA}{\mathcal{A}}
\newcommand{\bD}{\mathbf{D}}
\newcommand{\bX}{\mathbf{X}}
\newcommand{\bx}{\mathbf{x}}
\newcommand{\tX}{\widetilde{\mathbf{X}}}
\newcommand{\tx}{\widetilde{\mathbf{x}}}
\newcommand{\Sh}{\mathbb{S}}
\newcommand{\hphi}{\widehat{\varphi}}
\newcommand{\hS}{\widehat{S}}
\newcommand{\hsigma}{\widehat{\sigma}}
\newcommand{\hL}{\widehat{L}}
\newcommand{\hg}{\widehat{g}}
\newcommand{\hd}{\widehat{d}}
\newcommand{\hc}{\widehat{c}}
\newcommand{\Dn}{\mathcal{D}_n}
\newcommand{\eqdef}{\coloneqq}
\DeclareMathOperator*{\argmin}{arg\,min}
\DeclareMathOperator{\Card}{Card}
\DeclareMathOperator{\Var}{Var}
\DeclareMathOperator{\Id}{Id}

\begin{document}

% ===== début chapitres/00_titre.tex =====
\begin{titlepage}
\centering
{\large\scshape Aix-Marseille Université\par}
\vspace{0.2cm}
{\normalsize Master de Mathématiques\par}
\vspace{0.1cm}
{\normalsize Mémoire de Master 2\par}

\vspace{3cm}
\rule{\linewidth}{0.6pt}\par
\vspace{0.5cm}
{\LARGE\bfseries\color{bleuamu} Classification supervisée de données fonctionnelles par ondelettes\par}
\vspace{0.4cm}
{\large Règle de représentation commune fondée sur la fréquence d'utilisation\\ des caractéristiques temps-fréquence\par}
\vspace{0.5cm}
\rule{\linewidth}{0.6pt}\par

\vspace{2.5cm}
{\Large Hippolyte \textsc{Guigon}\par}

\vfill
\begin{tabular}{rl}
\textit{Encadrant :} & Prof. Jean-Marc \textsc{Freyermuth} \\[2pt]
\textit{Article de référence :} & A. Berlinet, G. Biau et L. Rouvière,\\
& \emph{Functional Supervised Classification with Wavelets}
\end{tabular}

\vspace{1.5cm}
{\normalsize Année universitaire 2025--2026\par}
\end{titlepage}
\cleardoublepage
% ===== fin chapitres/00_titre.tex =====


\pagenumbering{roman}
% ===== début chapitres/00_resume.tex =====
\chapter*{Résumé}
\addcontentsline{toc}{chapter}{Résumé}

Ce mémoire étudie la classification supervisée de courbes par la méthode de \citet{berlinet2008functional} : les courbes sont décomposées sur une base d'ondelettes, les fonctions de base sont réordonnées à l'aide de l'échantillon d'apprentissage, puis la dimension et la règle de classification sont choisies sur un échantillon de validation. Nous redémontrons en détail l'inégalité oracle et la consistance de cette procédure, en explicitant une hypothèse d'équivariance restée implicite. Nous formalisons l'emploi du log-périodogramme comme représentation spectrale des séries temporelles. Nous proposons enfin une nouvelle règle de représentation commune, fondée sur la fréquence à laquelle chaque fonction d'ondelette est utilisée par les courbes, après normalisation et avec un seuil global. Nous montrons qu'elle conserve les garanties de l'article et qu'elle possède en plus une concentration exponentielle sans hypothèse de moment, une influence bornée des courbes aberrantes et une invariance par changement d'échelle. Les expériences reproduisent les résultats de l'article et confirment ces propriétés : à performances comparables sur données propres, la nouvelle règle résiste à une contamination qui dégrade fortement la règle d'origine.

\paragraph{Mots-clés.} Données fonctionnelles, classification supervisée, ondelettes, périodogramme, minimisation du risque empirique, découpage de l'échantillon, robustesse.

\vspace{1cm}

\begin{otherlanguage}{english}
\section*{Abstract}

This thesis studies supervised classification of curves with the method of \citet{berlinet2008functional}: curves are expanded on a wavelet basis, the basis functions are reordered using the training sample, and the dimension and the classification rule are selected on a validation sample. We give detailed proofs of the oracle inequality and of the consistency of this procedure, and make explicit an equivariance assumption that was left implicit. We formalise the log-periodogram as a spectral representation of time series. Finally, we propose a new common-representation rule based on how often each wavelet is used by the curves, after normalisation and with a global threshold. We prove that it keeps the guarantees of the original method and additionally enjoys exponential concentration without moment assumptions, bounded influence of outlying curves and scale invariance. The experiments reproduce the results of the paper and confirm these properties: with comparable accuracy on clean data, the new rule withstands a contamination that severely degrades the original one.

\paragraph{Keywords.} Functional data, supervised classification, wavelets, periodogram, empirical risk minimisation, data splitting, robustness.
\end{otherlanguage}

\vspace{1cm}

\section*{Remerciements et usage d'outils}

Je remercie le Professeur Jean-Marc Freyermuth pour son encadrement et ses conseils tout au long de ce travail.

\noindent Un assistant d'intelligence artificielle (Claude, Anthropic) a été utilisé pour aider à l'écriture du code, à la mise en forme en \LaTeX{} et à la rédaction. Les choix méthodologiques ont été discutés avec l'encadrant ; les démonstrations et les résultats numériques ont été vérifiés et sont assumés par l'auteur.

\cleardoublepage
% ===== fin chapitres/00_resume.tex =====

\tableofcontents
\cleardoublepage
\pagenumbering{arabic}

% ===== début chapitres/01_introduction.tex =====
\chapter{Introduction}\label{chap-intro}

\section{Classer des courbes}

Dans de nombreux problèmes de classification, une observation n'est pas un vecteur de quelques mesures mais une courbe entière : un électrocardiogramme, un enregistrement de parole, un spectre de fréquences. On observe un échantillon $(X_1,Y_1),\dots,(X_N,Y_N)$ de copies indépendantes d'un couple $(X,Y)$, où $X$ prend ses valeurs dans un espace de Hilbert séparable de dimension infinie, typiquement $\Lp$, et $Y\in\{0,1\}$ est une étiquette. Le but est de construire, à partir de cet échantillon, une règle $g$ qui prédit l'étiquette d'une nouvelle courbe avec une probabilité d'erreur aussi proche que possible de l'erreur de Bayes $L^*$.

Ce cadre relève de l'analyse des données fonctionnelles \citep{ramsay2005functional,ferraty2006nonparametric}. La dimension infinie de l'espace des observations rend la plupart des méthodes classiques inopérantes telles quelles : les règles d'apprentissage souffrent du fléau de la dimension et certaines méthodes universellement consistantes dans $\R^d$ cessent de l'être dans un espace de Hilbert général \citep{abraham2006kernel,cerou2006nearest}. Une étape de réduction de dimension est donc indispensable avant toute classification.

\section{L'approche de Berlinet, Biau et Rouvière}

L'article de \citet{berlinet2008functional}, qui sert de point de départ à ce mémoire, propose une stratégie simple et justifiée théoriquement :
\begin{enumerate}
  \item chaque courbe est décomposée sur une base d'ondelettes orthonormée et tronquée à un niveau de résolution $J$, ce qui fournit $2^J$ coefficients ;
  \item les fonctions de base sont réordonnées selon un critère calculé sur l'échantillon d'apprentissage, l'\emph{énergie empirique} $\sum_i X_{ij}^2$ de chaque coefficient ;
  \item la dimension effective $d$ et la règle de classification $g$ sont choisies conjointement en minimisant l'erreur empirique sur un échantillon de validation indépendant.
\end{enumerate}
Toutes les courbes sont alors représentées par les mêmes $d$ coefficients : c'est la \emph{représentation commune}. Les auteurs démontrent une inégalité oracle pour la règle sélectionnée et en déduisent sa consistance.

\section{Objet du mémoire}

Ce mémoire a trois objectifs, fixés avec mon encadrant.

\paragraph{Reproduire la théorie de l'article.} Le \cref{chap-cadre} rassemble les outils de la reconnaissance automatique des formes \citep{devroye1988automatic,devroye1996probabilistic} et le \cref{chap-theorie} redémontre en détail l'inégalité oracle, le lemme d'approximation et le corollaire de consistance de \citet{berlinet2008functional}. Cette relecture fait apparaître une hypothèse implicite dans la preuve du corollaire, que nous explicitons (\cref{rem-equivariance}).

\paragraph{Travailler dans le domaine spectral.} Les données de parole de l'article sont des log-périodogrammes et non des signaux bruts. Nous formalisons cette étape d'ingénierie des caractéristiques à partir de \citet{shumway2011time} (\cref{sec-periodogramme}) et nous comparons, sur une vraie série temporelle (ECG200), le domaine temporel et le domaine spectral.

\paragraph{Proposer une nouvelle règle de représentation commune.} La règle de l'énergie n'est pas robuste : une seule courbe présentant un coefficient anormalement grand suffit à faire entrer l'indice correspondant dans la représentation commune. Nous proposons de classer les fonctions d'ondelettes selon leur \emph{fréquence d'utilisation} dans l'échantillon : après avoir ramené toutes les courbes à la même énergie, on compte la proportion de courbes dont le coefficient dépasse un \emph{seuil global} $\tau$. Nous montrons que cette règle hérite de toutes les garanties de l'article (\cref{thm-oracle-usage,cor-consistance-usage}) et qu'elle possède en plus des propriétés que la règle de l'énergie n'a pas : concentration exponentielle sans hypothèse de moment (\cref{thm-stabilite}), influence bornée d'une courbe aberrante (\cref{prop-robustesse}), invariance par changement d'échelle et parcimonie contrôlée (\cref{prop-proprietes-usage}).

\section{Plan}

Le \cref{chap-cadre} présente le cadre probabiliste et les inégalités de sélection par découpage de l'échantillon. Le \cref{chap-ondelettes} rappelle la construction des bases d'ondelettes sur $[0,1]$, la transformée discrète et le périodogramme. Le \cref{chap-methode} décrit la méthode de l'article et la nouvelle règle. Le \cref{chap-theorie} contient les résultats théoriques. Le \cref{chap-experiences} présente les expériences numériques : reproduction des résultats de l'article sur les données de parole et sur les données simulées, comparaison des domaines temporel et spectral, robustesse et comportement en fonction de la taille d'échantillon. L'annexe~\ref{chap-implementation} décrit l'implémentation.

\paragraph{Code et reproductibilité.} Tout le code est disponible sur le dépôt \url{https://github.com/HippolyteGuigon/functional_supervised_classification_with_wavelets}. Les commandes permettant de reproduire chaque tableau et chaque figure sont données en annexe.
% ===== fin chapitres/01_introduction.tex =====

% ===== début chapitres/02_cadre.tex =====
\chapter{Cadre : classification et sélection par validation}\label{chap-cadre}

Ce chapitre fixe le cadre probabiliste et démontre les trois inégalités sur lesquelles repose toute l'analyse : la réduction de la sélection à une déviation uniforme (\cref{lem-selection}), et deux bornes de cette déviation, pour une famille finie (\cref{prop-hoeffding}) et pour une famille de capacité finie au sens de Vapnik--Chervonenkis (\cref{prop-vc}). Nous suivons \citet[chapitres 2, 8, 12 et 13]{devroye1996probabilistic}.

\section{Risque et règle de Bayes}

Soit $\cF$ un espace de Hilbert séparable muni de sa tribu borélienne, et $(X,Y)$ un couple aléatoire à valeurs dans $\cF\times\{0,1\}$. On note
\[
  \eta(x) \eqdef \PP(Y=1\mid X=x), \qquad x\in\cF,
\]
une version de la probabilité conditionnelle. Une \emph{règle} est une application mesurable $g:\cF\to\{0,1\}$ ; sa probabilité d'erreur est $L(g)\eqdef\PP\{g(X)\neq Y\}$.

\begin{proposition}[Règle de Bayes]\label{prop-bayes}
La règle $g^*(x)=\1\{\eta(x)>1/2\}$ vérifie $L(g^*)\le L(g)$ pour toute règle $g$, et
\[
  L^* \eqdef \inf_{g} L(g) = L(g^*) = \E\bigl[\min\{\eta(X),1-\eta(X)\}\bigr].
\]
\end{proposition}

\begin{proof}
Pour une règle $g$ et $x\in\cF$, en conditionnant par $X=x$,
\[
  \PP\{g(X)\neq Y\mid X=x\} = \1\{g(x)=1\}\bigl(1-\eta(x)\bigr) + \1\{g(x)=0\}\,\eta(x) \;\ge\; \min\{\eta(x),1-\eta(x)\},
\]
avec égalité pour $g=g^*$. On intègre par rapport à la loi de $X$.
\end{proof}

La même preuve montre que, pour toute application mesurable $T:\cF\to\R^p$, l'erreur de Bayes du problème de prédiction de $Y$ à partir de $T(X)$ vaut
\begin{equation}\label{eq-bayes-T}
  L^*_{T} \eqdef \inf_{s:\R^p\to\{0,1\}}\PP\{s(T(X))\neq Y\} = \E\bigl[\min\{\eta_T,1-\eta_T\}\bigr],
  \qquad \eta_T\eqdef\PP\bigl(Y=1\mid T(X)\bigr),
\end{equation}
et que $L^*_T\ge L^*$, puisque toute règle de la forme $s\circ T$ est une règle sur $\cF$.

\section{Apprentissage automatique par découpage de l'échantillon}\label{sec-decoupage}

Suivant \citet{devroye1988automatic} et \citet{berlinet2008functional}, on dispose de $n+m$ copies indépendantes $(X_1,Y_1),\dots,(X_{n+m},Y_{n+m})$ de $(X,Y)$, découpées en une \emph{séquence d'apprentissage} $\Dn=\{(X_i,Y_i)\}_{i\le n}$ et une \emph{séquence de validation} $\{(X_i,Y_i)\}_{n<i\le n+m}$. On se donne une collection $\bD_n$, éventuellement infinie, de règles construites à partir de $\Dn$ seulement : chaque $g\in\bD_n$ est de la forme $x\mapsto g(x;\Dn)$. On note
\[
  L_n(g)\eqdef\PP\{g(X)\neq Y\mid \Dn\},
  \qquad
  \hL_{n,m}(g)\eqdef\frac1m\sum_{i=n+1}^{n+m}\1\{g(X_i)\neq Y_i\},
\]
où $(X,Y)$ est indépendant des données. La règle retenue $\hg$ minimise l'erreur de validation :
\begin{equation}\label{eq-erm}
  \hL_{n,m}(\hg) = \min_{g\in\bD_n}\hL_{n,m}(g).
\end{equation}
Sa probabilité d'erreur conditionnelle à toutes les données est $L_{n+m}(\hg)\eqdef\PP\{\hg(X)\neq Y\mid (X_i,Y_i)_{i\le n+m}\}$.

\begin{remarque}[Mesurabilité et existence du minimum]\label{rem-mesurabilite}
Lorsque $\bD_n$ est infinie, le minimum dans \eqref{eq-erm} est atteint car $\hL_{n,m}$ ne prend que $m+1$ valeurs. Nous supposons, comme \citet{devroye1996probabilistic}, que les suprema sur $\bD_n$ qui apparaissent ci-dessous sont mesurables ; c'est le cas dans toutes les situations considérées dans ce mémoire, où $\bD_n$ est finie ou paramétrée de façon continue par un ensemble séparable.
\end{remarque}

Le point clé est que, \emph{conditionnellement à $\Dn$}, les règles de $\bD_n$ sont des fonctions fixées et les couples de validation sont indépendants, de même loi que $(X,Y)$ et indépendants de $\Dn$. Pour $g\in\bD_n$ fixée, $m\hL_{n,m}(g)$ suit donc, conditionnellement à $\Dn$, une loi binomiale de paramètres $m$ et $L_n(g)$.

\begin{lemme}[{\citealp[lemme 8.2]{devroye1996probabilistic}}]\label{lem-selection}
Avec les notations précédentes,
\[
  L_{n+m}(\hg) - \inf_{g\in\bD_n}L_n(g) \;\le\; 2\sup_{g\in\bD_n}\bigl|\hL_{n,m}(g)-L_n(g)\bigr|.
\]
\end{lemme}

\begin{proof}
Notons $Z\eqdef\sup_{g\in\bD_n}|\hL_{n,m}(g)-L_n(g)|$. La règle $\hg$ est un élément (aléatoire) de $\bD_n$, et $L_{n+m}(\hg)$ est la valeur en $\hg$ de l'application $g\mapsto L_n(g)$ ; donc $L_{n+m}(\hg)\le \hL_{n,m}(\hg)+Z$. Pour toute $g\in\bD_n$, la définition \eqref{eq-erm} donne $\hL_{n,m}(\hg)\le \hL_{n,m}(g)\le L_n(g)+Z$. Ainsi $L_{n+m}(\hg)\le L_n(g)+2Z$ pour toute $g\in\bD_n$, et l'on passe à l'infimum en $g$.
\end{proof}

Il reste à contrôler l'espérance conditionnelle $\E_n Z$, où $\E_n$ désigne l'espérance conditionnelle à $\Dn$. Les deux bornes qui suivent reposent sur le même argument d'intégration.

\begin{lemme}\label{lem-integration}
Soit $Z\ge0$ une variable aléatoire telle que $\PP\{Z>\varepsilon\}\le A e^{-b\varepsilon^2}$ pour tout $\varepsilon>0$, avec $A\ge1$ et $b>0$. Alors, en posant $u=\sqrt{\log(A)/b}$,
\[
  \E Z \le u + \frac{1}{2bu} \qquad (u>0).
\]
\end{lemme}

\begin{proof}
Pour tout $u>0$, $\E Z=\int_0^\infty\PP\{Z>\varepsilon\}\,d\varepsilon\le u+\int_u^\infty Ae^{-b\varepsilon^2}d\varepsilon$. Pour $\varepsilon\ge u$, on a $e^{-b\varepsilon^2}\le\frac{\varepsilon}{u}e^{-b\varepsilon^2}$, d'où
$\int_u^\infty e^{-b\varepsilon^2}d\varepsilon\le\frac{1}{u}\int_u^\infty\varepsilon e^{-b\varepsilon^2}d\varepsilon=\frac{e^{-bu^2}}{2bu}$. Avec $u^2=\log(A)/b$, on a $Ae^{-bu^2}=1$.
\end{proof}

\begin{proposition}[Famille finie]\label{prop-hoeffding}
Si $\Card\bD_n\le N_n<\infty$, alors
\[
  \E_n\Bigl\{\sup_{g\in\bD_n}\bigl|\hL_{n,m}(g)-L_n(g)\bigr|\Bigr\}\le\sqrt{\frac{\log(2N_n)}{2m}}+\frac{1}{\sqrt{8m\log(2N_n)}}.
\]
\end{proposition}

\begin{proof}
Conditionnellement à $\Dn$, pour chaque $g$ fixée, $\hL_{n,m}(g)$ est une moyenne de $m$ variables de Bernoulli indépendantes d'espérance $L_n(g)$. L'inégalité de Hoeffding donne $\PP_n\{|\hL_{n,m}(g)-L_n(g)|>\varepsilon\}\le2e^{-2m\varepsilon^2}$, et une borne d'union sur les au plus $N_n$ règles donne $\PP_n\{Z>\varepsilon\}\le2N_ne^{-2m\varepsilon^2}$. Le \cref{lem-integration} avec $A=2N_n$ et $b=2m$ donne $u=\sqrt{\log(2N_n)/(2m)}$ et $\frac{1}{2bu}=\frac{1}{4m u}=\frac{1}{\sqrt{8m\log(2N_n)}}$.
\end{proof}

Lorsque $\bD_n$ est infinie (plus proches voisins avec tous les $k$, arbres de classification), on mesure sa richesse par le coefficient d'éclatement.

\begin{definition}[Coefficient d'éclatement]\label{def-shatter}
Pour une classe $\cA$ de parties mesurables d'un ensemble $\mathcal{Z}$ et un entier $N$,
\[
  \Sh_{\cA}(N) \eqdef \max_{z_1,\dots,z_N\in\mathcal{Z}}\Card\bigl\{\{z_1,\dots,z_N\}\cap A : A\in\cA\bigr\}.
\]
Pour une collection de règles $\bD_n$ fixée (c'est-à-dire conditionnellement à $\Dn$), on note $\cC_n\eqdef\{\{x:g(x)=1\}:g\in\bD_n\}$ et $\Sh_{\cC_n}(N)$ le coefficient correspondant.
\end{definition}

\begin{lemme}\label{lem-shatter-proprietes}
\begin{enumerate}[label=(\roman*)]
  \item Soit $\cA_n\eqdef\bigl\{\{(x,y):g(x)\neq y\}:g\in\bD_n\bigr\}$. Alors $\Sh_{\cA_n}(N)\le\Sh_{\cC_n}(N)$.
  \item Si $\bD_n=\bigcup_{\ell=1}^L\bD_n^{(\ell)}$, alors $\Sh_{\cC_n}(N)\le\sum_{\ell=1}^L\Sh_{\cC_n^{(\ell)}}(N)$.
  \item Soit $T:\cF\to\R^d$ mesurable et $\bD$ une classe de règles sur $\R^d$, de classe d'ensembles associée $\cC$. La classe $\{g\circ T:g\in\bD\}$ a une classe d'ensembles associée $\cC_T$ qui vérifie $\Sh_{\cC_T}(N)\le\Sh_{\cC}(N)$.
\end{enumerate}
\end{lemme}

\begin{proof}
(i) Fixons $(x_1,y_1),\dots,(x_N,y_N)$. La trace de l'ensemble d'erreur de $g$ est $\{(x_i,y_i):\1\{g(x_i)=1\}\neq y_i\}$ ; elle est entièrement déterminée par la trace $\{x_i:g(x_i)=1\}$ de $\cC_n$. L'application qui envoie la seconde trace sur la première est donc surjective de l'ensemble des traces de $\cC_n$ sur $\{x_1,\dots,x_N\}$ vers l'ensemble des traces de $\cA_n$, d'où l'inégalité.
(ii) Toute trace de $\cC_n$ est une trace d'au moins une des classes $\cC_n^{(\ell)}$.
(iii) La trace de $\{x:g(T(x))=1\}$ sur $\{x_1,\dots,x_N\}$ est l'image réciproque par $T$ de la trace de $\{u:g(u)=1\}$ sur $\{T(x_1),\dots,T(x_N)\}$, ensemble d'au plus $N$ points ; deux règles ayant la même trace sur les $T(x_i)$ ont la même trace sur les $x_i$.
\end{proof}

\begin{theoreme}[Inégalité de Vapnik--Chervonenkis, {\citealp[théorème 12.5]{devroye1996probabilistic}}]\label{thm-vc}
Soit $Z_1,\dots,Z_m$ des variables i.i.d. de loi $\nu$ sur $\mathcal{Z}$, $\nu_m$ leur mesure empirique, et $\cA$ une classe de parties mesurables. Pour tout $\varepsilon>0$,
\[
  \PP\Bigl\{\sup_{A\in\cA}|\nu_m(A)-\nu(A)|>\varepsilon\Bigr\}\le 4\,\Sh_{\cA}(2m)\,e^{-m\varepsilon^2/8}.
\]
\end{theoreme}

Nous admettons ce résultat classique, dont la preuve repose sur une double symétrisation et sur l'inégalité de Hoeffding \citep{vapnik1971uniform,devroye1996probabilistic}.

\begin{proposition}[Famille de capacité finie]\label{prop-vc}
\[
  \E_n\Bigl\{\sup_{g\in\bD_n}\bigl|\hL_{n,m}(g)-L_n(g)\bigr|\Bigr\}\le\sqrt{\frac{8\log\bigl(4\Sh_{\cC_n}(2m)\bigr)}{m}}+\frac{1}{\sqrt{(m/2)\log\bigl(4\Sh_{\cC_n}(2m)\bigr)}}.
\]
\end{proposition}

\begin{proof}
Conditionnellement à $\Dn$, on applique le \cref{thm-vc} aux couples de validation $Z_i=(X_{n+i},Y_{n+i})$ et à la classe $\cA_n$ du \cref{lem-shatter-proprietes}\,(i) : $\nu(\{(x,y):g(x)\neq y\})=L_n(g)$ et $\nu_m(\cdot)=\hL_{n,m}(g)$. Avec (i), $\PP_n\{Z>\varepsilon\}\le4\Sh_{\cC_n}(2m)e^{-m\varepsilon^2/8}$. Le \cref{lem-integration} avec $A=4\Sh_{\cC_n}(2m)\ge4$ et $b=m/8$ donne $u=\sqrt{8\log(4\Sh_{\cC_n}(2m))/m}$ et
$\frac{1}{2bu}=\frac{4}{mu}=\frac{1}{\sqrt{(m/2)\log(4\Sh_{\cC_n}(2m))}}$.
\end{proof}

\begin{lemme}\label{lem-monotonie}
Pour $m\ge1$, la fonction $f_m(u)\eqdef\sqrt{8\log(4u)/m}+1/\sqrt{(m/2)\log(4u)}$ est croissante sur $[1,\infty[$.
\end{lemme}

\begin{proof}
Avec $\ell=\sqrt{\log(4u)}\ge\sqrt{\log4}$, on a $f_m=\sqrt{2/m}\,(2\ell+1/\ell)$, et $\ell\mapsto2\ell+1/\ell$ est croissante pour $\ell^2>1/2$, ce qui est le cas car $\log4\approx1{,}386$. Enfin $u\mapsto\ell$ est croissante.
\end{proof}

Ce lemme permet de remplacer $\Sh_{\cC_n}(2m)$ par n'importe quel majorant dans la \cref{prop-vc}.
% ===== fin chapitres/02_cadre.tex =====

% ===== début chapitres/03_ondelettes.tex =====
\chapter{Ondelettes et périodogramme}\label{chap-ondelettes}

Ce chapitre rassemble les deux transformations appliquées aux courbes avant la classification : la décomposition sur une base d'ondelettes orthonormée, qui fournit les coefficients $X_{ij}$, et le périodogramme, qui fait passer une série temporelle dans le domaine des fréquences. Nous insistons sur une propriété qui joue un rôle central dans la suite : la transformée utilisée doit être \emph{orthogonale}, afin que l'énergie d'une courbe soit exactement la somme des carrés de ses coefficients.

\section{Analyse multirésolution et bases d'ondelettes}\label{sec-amr}

\begin{definition}[Analyse multirésolution]
Une analyse multirésolution de $L^2(\R)$ est une suite croissante de sous-espaces fermés $(V_j)_{j\in\mathbb{Z}}$ telle que
\begin{enumerate}[label=(\roman*)]
  \item $\bigcap_j V_j=\{0\}$ et $\bigcup_j V_j$ est dense dans $L^2(\R)$ ;
  \item $f\in V_j\iff f(2\,\cdot)\in V_{j+1}$ ;
  \item il existe une fonction d'échelle $\phi\in V_0$ telle que $\{\phi(\cdot-k)\}_{k\in\mathbb{Z}}$ soit une base orthonormée de $V_0$.
\end{enumerate}
\end{definition}

On note $W_j$ le supplémentaire orthogonal de $V_j$ dans $V_{j+1}$. Il existe une \emph{ondelette mère} $\psi$ telle que $\{\psi_{j,k}\}_{k\in\mathbb{Z}}$, avec $\psi_{j,k}(t)=2^{j/2}\psi(2^jt-k)$, soit une base orthonormée de $W_j$ \citep{mallat1989theory,daubechies1992ten}. Les ondelettes de Daubechies \texttt{db}$p$ sont à support compact de longueur $2p-1$ et possèdent $p$ moments nuls ; \texttt{db1} est la base de Haar \citep{haar1910theorie}.

\paragraph{Bases sur $[0,1]$ par périodisation.} Pour une fonction $h$ de $L^2(\R)$, on pose $h^{\mathrm{per}}(t)=\sum_{\ell\in\mathbb{Z}}h(t+\ell)$. Les familles
\[
  \bigl\{\phi_{0,0}^{\mathrm{per}}\bigr\}\;\cup\;\bigl\{\psi_{j,k}^{\mathrm{per}}:\ j\ge0,\ 0\le k<2^j\bigr\}
\]
forment une base orthonormée de $\Lp$ \citep[section 9.3]{daubechies1992ten}, où $\phi^{\mathrm{per}}_{0,0}\equiv1$. C'est le cadre de \citet{berlinet2008functional} : toute courbe $X\in\Lp$ s'écrit
\begin{equation}\label{eq-decomposition}
  X = \eta\,\phi_{0,0} + \sum_{j\ge0}\sum_{k=0}^{2^j-1}\zeta_{j,k}\psi_{j,k},
  \qquad \zeta_{j,k}=\langle X,\psi_{j,k}\rangle,\quad \eta=\langle X,\phi_{0,0}\rangle,
\end{equation}
la série convergeant dans $\Lp$. On réindexe la base en
\[
  \{\psi_1,\psi_2,\psi_3,\dots\}=\{\phi_{0,0},\psi_{0,0},\psi_{1,0},\psi_{1,1},\psi_{2,0},\dots\},
\]
de sorte que $X=\sum_{j\ge1}X_j\psi_j$ avec $X_j=\langle X,\psi_j\rangle$, et que les $2^J$ premières fonctions engendrent exactement $V_J$.

\begin{proposition}[Parseval et projection]\label{prop-parseval}
Pour tout $X\in\Lp$ et tout $J\ge0$, notons $P_J X=\sum_{j=1}^{2^J}X_j\psi_j$ la projection orthogonale de $X$ sur $V_J$. Alors
\[
  \|P_JX\|^2=\sum_{j=1}^{2^J}X_j^2,\qquad \|X-P_JX\|^2=\sum_{j>2^J}X_j^2\xrightarrow[J\to\infty]{}0.
\]
\end{proposition}

\begin{proof}
Conséquence directe de l'orthonormalité de la base et de la convergence dans $\Lp$ de \eqref{eq-decomposition}.
\end{proof}

\section{Transformée discrète et orthogonalité}\label{sec-dwt}

En pratique, une courbe est observée en $T$ points équidistants, $x=(x(t_1),\dots,x(t_T))$. L'algorithme pyramidal de \citet{mallat1989theory} calcule les coefficients par filtrages successifs à l'aide des filtres $h$ (passe-bas) et $g$ (passe-haut) associés à $\phi$ et $\psi$ : à chaque niveau, les coefficients d'approximation sont convolués par $h$ et $g$ puis sous-échantillonnés d'un facteur 2.

Le traitement des bords est crucial. Avec l'extension \emph{périodique}, un pas de l'algorithme sur un vecteur de longueur paire $N$ s'écrit $a\mapsto(Ha,Ga)$, où $H,G\in\R^{N/2\times N}$ sont des matrices circulantes sous-échantillonnées vérifiant $HH^\top=GG^\top=I$ et $HG^\top=0$ (conditions de filtres miroirs en quadrature). La matrice $\binom{H}{G}$ est donc orthogonale. Tant que la longueur reste paire, on peut itérer.

\begin{proposition}\label{prop-dwt-orthogonale}
Soit $T=2^{L}q$ avec $q$ impair. La transformée en ondelettes discrète périodisée à $L'\le L$ niveaux est une application linéaire orthogonale $W:\R^T\to\R^T$. En particulier $\|Wx\|_2=\|x\|_2$ pour tout $x\in\R^T$.
\end{proposition}

\begin{proof}
Chaque niveau $\ell\le L'$ agit sur un vecteur d'approximation de longueur $T/2^{\ell-1}$, qui est paire car $\ell-1<L$ ; il est orthogonal d'après ce qui précède. On le complète par l'identité sur les coefficients de détail déjà calculés. $W$ est la composée de ces applications orthogonales.
\end{proof}

\begin{remarque}[Conséquence pratique]\label{rem-mode-symmetric}
L'extension par défaut de la bibliothèque PyWavelets (mode \texttt{symmetric}) prolonge le signal par symétrie avant chaque filtrage et renvoie \emph{plus} de $T$ coefficients. La transformée n'est alors plus orthogonale. Sur un bruit blanc de longueur 96 (taille des ECG200), l'énergie des coefficients \texttt{db4} dépasse celle du signal d'environ 25\,\% ; pour \texttt{db10} elle est multipliée par 2,7. Les coefficients de bord sont artificiellement surpondérés, ce qui fausse le classement par énergie \eqref{eq-energie} et la normalisation des courbes utilisée au \cref{chap-methode}. Toutes les expériences de ce mémoire utilisent donc le mode \texttt{periodization} avec $L=\max\{\ell:2^\ell\mid T\}$ niveaux ($L=5$ pour $T=96$, $L=8$ pour $T=256$, $L=10$ pour $T=1024$) ; la propriété de Parseval est vérifiée numériquement par les tests du dépôt.
\end{remarque}

Dans la suite on identifie une courbe à son vecteur de coefficients $\bX=(X_1,\dots,X_p)\in\R^p$, $p=2^J$ dans le cadre théorique et $p=T$ dans les expériences.

\section{Le périodogramme}\label{sec-periodogramme}

Les données de parole de \citet{berlinet2008functional}, reprises de \citet{hastie1995penalized}, ne sont pas des enregistrements bruts mais des \emph{log-périodogrammes}. Nous décrivons ici cette transformation, qui constitue une étape d'ingénierie des caractéristiques, en suivant \citet[section 4.3]{shumway2011time}.

\begin{definition}[Transformée de Fourier discrète et périodogramme]
Soit $x_1,\dots,x_T$ une série observée. Aux fréquences de Fourier $\omega_k=k/T$, $k=0,\dots,T-1$, on définit
\[
  d(\omega_k)=T^{-1/2}\sum_{t=1}^{T}x_t\,e^{-2\pi i\omega_kt},
  \qquad
  I(\omega_k)=|d(\omega_k)|^2 .
\]
\end{definition}

Comme $d(0)=\sqrt T\,\bar x$, on a $I(0)=T\bar x^2$ ; pour $k\ne0$, $\sum_te^{-2\pi i\omega_kt}=0$, donc $I(\omega_k)$ ne change pas si l'on retranche la moyenne. Le périodogramme décompose la variance empirique selon les fréquences :
\[
  \sum_{t=1}^T(x_t-\bar x)^2=\sum_{k=1}^{T-1}I(\omega_k),
\]
par la formule de Parseval pour la transformée de Fourier discrète.

\begin{proposition}[{\citealp[propriété 4.5]{shumway2011time}}]\label{prop-periodogramme}
Soit $(x_t)$ un processus linéaire $x_t=\sum_{j\in\mathbb Z}\psi_jw_{t-j}$ avec $(w_t)$ i.i.d. centré de variance $\sigma_w^2$, $\sum_j|\psi_j|<\infty$, de fonction d'autocovariance $\gamma$ vérifiant $\sum_h|h|\,|\gamma(h)|<\infty$ et de densité spectrale $f$. Pour des fréquences distinctes fixées $0<\nu_1,\dots,\nu_M<1/2$ telles que $f(\nu_\ell)>0$, et des fréquences de Fourier $\omega_{k_T(\ell)}\to\nu_\ell$,
\[
  \frac{2I(\omega_{k_T(\ell)})}{f(\nu_\ell)}\xrightarrow[T\to\infty]{\mathcal L}\chi^2_2,
\]
les $M$ limites étant indépendantes.
\end{proposition}

Deux conséquences guident notre usage du périodogramme.

\paragraph{Un estimateur non consistant.} On a $\E I(\omega_k)\to f(\nu)$ mais $\Var I(\omega_k)\to f(\nu)^2$ : la variance ne tend pas vers zéro. Le périodogramme est un estimateur sans biais asymptotique mais non consistant de la densité spectrale. Pour la classification, ce n'est pas un défaut rédhibitoire : on ne cherche pas à estimer $f$, mais à disposer d'une représentation dans laquelle les classes se séparent. La décomposition en ondelettes suivie d'une sélection de coefficients joue en partie le rôle du lissage qui rendrait l'estimateur consistant.

\paragraph{Stabilisation de la variance par le logarithme.} Si $2I/f\sim\chi^2_2$, alors $I/f$ suit une loi exponentielle $\mathcal E(1)$ et
\[
  \log I(\omega_k)=\log f(\omega_k)+\varepsilon_k,
  \qquad \E\varepsilon_k=-\gamma,\quad \Var\varepsilon_k=\frac{\pi^2}{6},
\]
où $\gamma\approx0{,}5772$ est la constante d'Euler. Le log-périodogramme est donc, asymptotiquement, le log de la densité spectrale plus un bruit additif de variance constante : c'est le cadre naturel d'un débruitage par ondelettes \citep{wahba1980automatic,moulin1994wavelet}. Nous travaillons toujours avec $\log(I+\epsilon)$, $\epsilon=10^{-12}$.

\paragraph{Conventions.} Nous retirons la moyenne de la série avant la transformée (retrait de tendance constant), puis nous supprimons l'ordonnée $I(0)$, nulle après centrage et qui écraserait l'échelle logarithmique. Pour une série de longueur $T$, on conserve les fréquences $\omega_k$, $1\le k\le\lfloor T/2\rfloor$, soit $48$ ordonnées pour les ECG200.

\paragraph{Invariance par translation.} Si $y_t=x_{t-s}$, les indices étant pris modulo $T$ (translation circulaire), alors
\[
  d_y(\omega_k)=e^{-2\pi i\omega_ks}\,d_x(\omega_k)\qquad\text{et donc}\qquad I_y=I_x .
\] Le périodogramme oublie la phase, donc la position temporelle des motifs. C'est un avantage lorsque l'information discriminante est fréquentielle (phonèmes), et un inconvénient lorsque la forme et la position d'un motif comptent (complexe QRS d'un électrocardiogramme). Nos expériences du \cref{sec-exp-ecg} confirment ce point.

\begin{remarque}[Données déjà spectrales]\label{rem-phoneme-spectral}
Les courbes du jeu de données \emph{phoneme} sont déjà des log-périodogrammes. Leur appliquer à nouveau la transformation produirait le périodogramme d'un log-périodogramme, objet sans interprétation physique ; dans une première version de nos expériences, cette erreur faisait chuter l'exactitude de $0{,}86$ à environ $0{,}6$. Ces données sont donc utilisées telles quelles : elles \emph{sont} la représentation spectrale.
\end{remarque}
% ===== fin chapitres/03_ondelettes.tex =====

% ===== début chapitres/04_methode.tex =====
\chapter{Représentation commune et sélection de modèle}\label{chap-methode}

\section{Schéma général}\label{sec-schema}

On fixe un niveau de résolution maximal $J$, éventuellement fonction de $n$, et l'on pose $p=2^J$. Chaque courbe $X_i$ est représentée par le vecteur de ses $p$ premiers coefficients
\[
  \bX_i=(X_{i1},\dots,X_{ip})\in\R^p,\qquad X_{ij}=\langle X_i,\psi_j\rangle,
\]
c'est-à-dire par sa projection $P_JX_i$ sur $V_J$ (\cref{prop-parseval}). Pour une permutation $\sigma$ de $\{1,\dots,p\}$ et $1\le d\le p$, on note
\[
  \Pi_{\sigma,d}:\R^p\to\R^d,\qquad \Pi_{\sigma,d}(\bx)=(x_{\sigma(1)},\dots,x_{\sigma(d)})
\]
la projection sur les $d$ coordonnées de rang le plus élevé selon $\sigma$. La méthode comporte trois étapes.

\begin{enumerate}
  \item \textbf{Familles de règles.} Pour chaque $d\le p$, on se donne une collection $\cD_n^{(d)}$ de règles de classification opérant dans $\R^d$ et construites à partir de $n$ points d'apprentissage de $\R^d\times\{0,1\}$ : $k$ plus proches voisins pour tous les $k$ (W-NN), analyse discriminante quadratique (W-QDA), arbres de classification (W-CART).
  \item \textbf{Réordonnancement.} La séquence d'apprentissage détermine une permutation $\hsigma$ (ou une famille de permutations $\hsigma_c$ indexée par un paramètre $c$) qui classe les fonctions de base de la plus à la moins informative. L'ensemble $\hS_d=\{\hsigma(1),\dots,\hsigma(d)\}$ est la \emph{représentation commune} en dimension $d$ : \emph{toutes} les courbes sont représentées par les mêmes $d$ coefficients.
  \item \textbf{Sélection.} La dimension, le paramètre et la règle sont choisis en minimisant l'erreur sur la séquence de validation :
  \begin{equation}\label{eq-selection}
    (\hc,\hd,\hg)\in\argmin_{c\in\Gamma,\ 1\le d\le p,\ g\in\cD_n^{(d)}}\ \frac1m\sum_{i=n+1}^{n+m}\1\bigl\{g\bigl(\Pi_{\hsigma_c,d}(\bX_i)\bigr)\neq Y_i\bigr\}.
  \end{equation}
\end{enumerate}

Dans l'étape 3, la règle $g\in\cD_n^{(d)}$ est entraînée sur les points $\bigl(\Pi_{\hsigma_c,d}(\bX_i),Y_i\bigr)_{i\le n}$. Conditionnellement à $\Dn$, la collection
\begin{equation}\label{eq-collection}
  \bD_n\eqdef\bigcup_{c\in\Gamma}\ \bigcup_{d=1}^{p}\ \bigl\{x\mapsto g\bigl(\Pi_{\hsigma_c,d}(\bX)\bigr):\ g\in\cD_n^{(d)}\bigr\}
\end{equation}
est donc une collection de règles fixées, et \eqref{eq-selection} est exactement la minimisation \eqref{eq-erm} du \cref{chap-cadre}. Il reste à spécifier le réordonnancement.

\section{La règle de l'énergie}\label{sec-energie}

\citet{berlinet2008functional} classent les fonctions de base selon leur énergie empirique :
\begin{equation}\label{eq-energie}
  \sum_{i=1}^nX_{i\hsigma(1)}^2\ \ge\ \sum_{i=1}^nX_{i\hsigma(2)}^2\ \ge\ \dots\ \ge\ \sum_{i=1}^nX_{i\hsigma(p)}^2 .
\end{equation}
Il n'y a pas de paramètre ($\Gamma$ est réduit à un point). Cette étape s'apparente au seuillage des coefficients d'ondelettes \citep{donoho1995wavelet}, avec une différence essentielle soulignée par les auteurs : le seuillage est \emph{global}, commun à toutes les courbes, alors qu'un seuillage individuel produirait une base significative différente pour chaque courbe et donc pas de représentation commune.

\paragraph{Fragilité.} Le critère \eqref{eq-energie} est une somme de carrés : il est dominé par les plus grandes valeurs. Une seule courbe aberrante, par exemple un artefact d'enregistrement localisé, possède quelques coefficients fins de très grande amplitude, et ces indices passent en tête du classement même si aucune autre courbe ne les utilise. De même, une courbe d'amplitude globale très supérieure aux autres impose sa propre base significative à tout l'échantillon. Enfin, la convergence de $\frac1n\sum_iX_{ij}^2$ vers $\E X_j^2$ requiert que ce moment soit fini ; sans hypothèse supplémentaire, aucune vitesse n'est disponible. Ces trois défauts sont formalisés et quantifiés au \cref{sec-proprietes}.

\section{Une règle fondée sur la fréquence d'utilisation}\label{sec-usage}

L'idée est de remplacer une mesure d'\emph{intensité} (combien d'énergie porte l'indice $j$ dans l'échantillon) par une mesure de \emph{fréquence} (dans combien de courbes l'indice $j$ est significatif). Pour que « significatif » ait le même sens pour toutes les courbes, on les ramène d'abord à la même énergie, puis on applique un seuil unique.

\begin{definition}[Normalisation]\label{def-normalisation}
Pour $\bx\in\R^p$, on pose $\tx\eqdef\bx/\|\bx\|$ si $\bx\neq0$ et $\tx\eqdef0$ sinon, où $\|\cdot\|$ est la norme euclidienne. D'après la \cref{prop-parseval}, $\|\bX_i\|=\|P_JX_i\|_{\Lp}$ : normaliser le vecteur des coefficients revient à normaliser la courbe projetée dans $\Lp$.
\end{definition}

\begin{definition}[Fréquence d'utilisation et règle d'usage]\label{def-usage}
Soit $c>0$. On dit que la courbe $x$ \emph{utilise} la fonction de base $\psi_j$ au niveau $c$ si
\begin{equation}\label{eq-indicatrice}
  s_j(\bx;c)\eqdef\1\Bigl\{|\tx_j|>\frac{c}{\sqrt p}\Bigr\}=1 .
\end{equation}
La \emph{fréquence d'utilisation empirique} de $\psi_j$ est
\begin{equation}\label{eq-usage}
  \hphi_j(c)\eqdef\frac1n\sum_{i=1}^n s_j(\bX_i;c)\in[0,1],
\end{equation}
et la permutation $\hsigma_c$ classe les indices par fréquence décroissante, les ex æquo étant départagés par l'ordre naturel (du plus grossier au plus fin) :
\[
  \hphi_{\hsigma_c(1)}(c)\ge\hphi_{\hsigma_c(2)}(c)\ge\dots\ge\hphi_{\hsigma_c(p)}(c).
\]
La représentation commune en dimension $d$ est $\hS_d(c)=\{\hsigma_c(1),\dots,\hsigma_c(d)\}$.
\end{definition}

\paragraph{Interprétation du seuil.} Comme $\|\tx\|=1$, la condition $|\tx_j|>c/\sqrt p$ équivaut à
\[
  x_j^2>\frac{c^2}{p}\,\|\bx\|^2 :
\]
la courbe utilise $\psi_j$ lorsque ce coefficient porte plus de $c^2$ fois la part d'énergie qu'il recevrait si l'énergie de la courbe était répartie uniformément sur les $p$ coefficients. Le paramètre $c$ est donc sans dimension et comparable d'un jeu de données à l'autre. Pour $c\ge\sqrt p$, aucun coefficient n'est jamais utilisé ; pour $c\to0$, tout coefficient non nul l'est. On prend donc $\Gamma\subset\,]0,\sqrt p[$ fini ; dans nos expériences, $\Gamma=\{0{,}5;\,1;\,1{,}5;\,2;\,3\}$, puis une grille élargie jusqu'à $c=8$ (\cref{sec-exp-grille}).

\paragraph{Seuil global ou seuil individuel.} Sur les courbes normalisées, le seuil $\tau=c/\sqrt p$ est le même pour toutes les courbes. Sur les courbes d'origine, il correspond au seuil $\tau_i=c\,\|\bX_i\|/\sqrt p$, proportionnel à la norme de chaque courbe. Une première version de ce travail utilisait des seuils individuels non liés à la norme (garder les $m$ plus grands coefficients de chaque courbe, ou un seuil universel de \citealt{donoho1994ideal}) ; le choix d'un seuil global après normalisation, retenu avec l'encadrant, simplifie à la fois l'interprétation et l'analyse théorique.

\paragraph{Ce que la règle ne fait pas.} La règle d'usage est \emph{non supervisée} : elle n'utilise pas les étiquettes. Nous avions aussi étudié une variante qui classe les indices par l'écart de fréquence d'utilisation entre classes, $|\hphi^{(1)}_j-\hphi^{(0)}_j|$. Elle donnait parfois de meilleures erreurs de prédiction, mais elle ne définit plus une représentation \emph{commune} décrivant les courbes : elle optimise directement la séparation, au risque de sur-ajuster le critère de sélection. Conformément à l'objectif fixé, qui est de choisir une représentation adaptée aux données même si elle n'est pas la plus performante en prédiction, nous ne la retenons pas.

\paragraph{Normalisation et classification.} La normalisation sert uniquement à \emph{choisir} les indices. Les règles $g\in\cD_n^{(d)}$ reçoivent les coefficients d'origine $\Pi_{\hsigma_c,d}(\bX)$, non normalisés. Ce choix est essentiel pour la théorie : la \cref{prop-normalisation-perte} montre que classer les courbes normalisées peut augmenter strictement l'erreur de Bayes.

\section{Récapitulatif de l'algorithme}

\begin{enumerate}[label=\arabic*.]
  \item Si les courbes sont des séries temporelles et que l'on travaille dans le domaine spectral, remplacer chaque courbe par son log-périodogramme (\cref{sec-periodogramme}).
  \item Calculer les coefficients d'ondelettes par la transformée périodisée orthogonale (\cref{prop-dwt-orthogonale}).
  \item Sur la séquence d'apprentissage seulement : normaliser les vecteurs de coefficients, calculer $\hphi_j(c)$ pour chaque $c\in\Gamma$ et en déduire $\hsigma_c$.
  \item Pour chaque $c\in\Gamma$, $d\le d_{\max}$ et $g\in\cD_n^{(d)}$ : entraîner $g$ sur $(\Pi_{\hsigma_c,d}(\bX_i),Y_i)_{i\le n}$ et calculer son erreur sur la séquence de validation.
  \item Retenir le triplet d'erreur de validation minimale (en cas d'égalité, la plus petite dimension).
\end{enumerate}

\begin{remarque}[Validation croisée]
Le code du dépôt propose aussi une sélection par validation croisée à $K$ blocs, où le classement est recalculé à l'intérieur de chaque bloc à partir de ses seules données d'apprentissage. C'est l'usage courant en pratique, mais les garanties du \cref{chap-theorie} portent sur le découpage apprentissage/validation, qui est aussi le protocole de \citet{berlinet2008functional}. Toutes les expériences du \cref{chap-experiences} utilisent donc ce découpage.
\end{remarque}
% ===== fin chapitres/04_methode.tex =====

% ===== début chapitres/05_theorie.tex =====
\chapter{Résultats théoriques}\label{chap-theorie}

Ce chapitre contient les résultats théoriques du mémoire. La \cref{sec-reproduction} redémontre les résultats de \citet{berlinet2008functional} dans un cadre qui couvre directement toute règle de réordonnancement calculée sur la séquence d'apprentissage, et en particulier la règle d'usage. La \cref{sec-proprietes} établit les propriétés propres à la règle d'usage, que la règle de l'énergie ne possède pas.

\section{Notations et hypothèses}\label{sec-notations}

On reprend les notations des \cref{chap-cadre,chap-methode}. $(X,Y)$ est à valeurs dans $\Lp\times\{0,1\}$, $\{\psi_j\}_{j\ge1}$ est la base d'ondelettes réindexée, $p=2^J$ et
\[
  T_J:\Lp\to\R^p,\qquad T_J(x)=\bigl(\langle x,\psi_1\rangle,\dots,\langle x,\psi_p\rangle\bigr),\qquad \bX=T_J(X).
\]
On note $L^*_p\eqdef L^*_{T_J}$ l'erreur de Bayes lorsque l'on observe seulement $\bX$ (\cref{eq-bayes-T}). Pour une permutation $\sigma$ de $\{1,\dots,p\}$, on écrit $\Pi_\sigma\eqdef\Pi_{\sigma,p}$ ; c'est une bijection linéaire de $\R^p$.

\begin{lemme}\label{lem-bayes-permutation}
Pour toute permutation $\sigma$ déterministe, l'erreur de Bayes de la prédiction de $Y$ à partir de $\Pi_\sigma(\bX)$ est égale à $L^*_p$.
\end{lemme}

\begin{proof}
Les règles $s:\R^p\to\{0,1\}$ appliquées à $\Pi_\sigma(\bX)$ sont exactement les règles $s\circ\Pi_\sigma$ appliquées à $\bX$, et $s\mapsto s\circ\Pi_\sigma$ est une bijection de l'ensemble des règles mesurables sur $\R^p$.
\end{proof}

\paragraph{Le cadre de sélection.} Un \emph{réordonnancement} est une famille $(\hsigma_c)_{c\in\Gamma}$ de permutations de $\{1,\dots,p\}$, indexée par un ensemble fini $\Gamma$ et fonction mesurable de $\Dn$ seulement. La règle de l'énergie correspond à $\Card\Gamma=1$ ; la règle d'usage à $\hsigma_c$ défini par la \cref{def-usage}. On note $\cC_n^{(d,c)}$ la classe des ensembles $\{u\in\R^d:g(u)=1\}$, $g\in\cD_n^{(d)}$, lorsque les règles de $\cD_n^{(d)}$ sont entraînées sur $(\Pi_{\hsigma_c,d}(\bX_i),Y_i)_{i\le n}$, et
\[
  \Sh^{(J)}_n(N)\eqdef\max_{c\in\Gamma}\ \sum_{d=1}^p\Sh_{\cC_n^{(d,c)}}(N).
\]
Dans les exemples usuels, $\Sh_{\cC_n^{(d,c)}}(N)$ admet un majorant qui ne dépend ni des données ni de $c$ (\cref{ex-nn,ex-arbres}).

\section{Inégalité oracle et consistance}\label{sec-reproduction}

\subsection{Inégalité oracle}

\begin{theoreme}[Inégalité oracle ; {\citealp[théorème 2.1]{berlinet2008functional}}, étendu]\label{thm-oracle}
Pour tout réordonnancement $(\hsigma_c)_{c\in\Gamma}$ fonction de $\Dn$ seulement, la règle $\hg$ définie par \eqref{eq-selection} vérifie
\begin{multline*}
  \E\{L_{n+m}(\hg)\}-L^*\ \le\ \underbrace{\bigl(L^*_p-L^*\bigr)}_{\text{approximation}}
  +\underbrace{\Bigl(\E\Bigl\{\inf_{g\in\bD_n}L_n(g)\Bigr\}-L^*_p\Bigr)}_{\text{estimation}}\\
  +\underbrace{2\,\E\biggl\{\sqrt{\frac{8\log\bigl(4\Card(\Gamma)\,\Sh^{(J)}_n(2m)\bigr)}{m}}+\frac{1}{\sqrt{(m/2)\log\bigl(4\Card(\Gamma)\,\Sh^{(J)}_n(2m)\bigr)}}\biggr\}}_{\text{sélection}},
\end{multline*}
où $\bD_n$ est la collection \eqref{eq-collection}.
\end{theoreme}

\begin{proof}
\emph{Étape 1 : capacité de $\bD_n$.} Fixons $\Dn$. Pour $c\in\Gamma$ et $d\le p$, la sous-collection $\{g\circ\Pi_{\hsigma_c,d}\circ T_J:g\in\cD_n^{(d)}\}$ est de la forme $\{g\circ T:g\in\cD_n^{(d)}\}$ avec $T=\Pi_{\hsigma_c,d}\circ T_J$ mesurable et fixé. Par le \cref{lem-shatter-proprietes}\,(iii), son coefficient d'éclatement est majoré par $\Sh_{\cC_n^{(d,c)}}$. Par le \cref{lem-shatter-proprietes}\,(ii) appliqué à la réunion \eqref{eq-collection},
\[
  \Sh_{\cC_n}(2m)\le\sum_{c\in\Gamma}\sum_{d=1}^p\Sh_{\cC_n^{(d,c)}}(2m)\le\Card(\Gamma)\,\Sh^{(J)}_n(2m).
\]

\emph{Étape 2 : sélection.} Conditionnellement à $\Dn$, $\bD_n$ est une collection de règles fixées et $\hg$ minimise l'erreur de validation sur $\bD_n$. Le \cref{lem-selection}, la \cref{prop-vc} et la croissance de $f_m$ (\cref{lem-monotonie}) donnent
\[
  \E_n\{L_{n+m}(\hg)\}-\inf_{g\in\bD_n}L_n(g)\le2f_m\bigl(\Card(\Gamma)\,\Sh^{(J)}_n(2m)\bigr).
\]
On prend l'espérance par rapport à $\Dn$.

\emph{Étape 3 : décomposition.} On écrit
$\E L_{n+m}(\hg)-L^*=(L^*_p-L^*)+\bigl(\E\inf_{\bD_n}L_n-L^*_p\bigr)+\bigl(\E L_{n+m}(\hg)-\E\inf_{\bD_n}L_n\bigr)$ et l'on majore le dernier terme par l'étape 2.
\end{proof}

\begin{remarque}\label{rem-oracle-lecture}
Trois points méritent d'être soulignés.
\begin{enumerate}[label=(\alph*)]
  \item Le terme d'approximation $L^*_p-L^*$ ne dépend \emph{pas} du réordonnancement. C'est la perte due à la troncature au niveau $J$.
  \item Le coût du paramètre $c$ est un terme en $\log\Card(\Gamma)$ sous la racine : choisir le seuil parmi une grille finie ne coûte presque rien.
  \item Le théorème ne dit pas quel réordonnancement est meilleur. L'intérêt d'un bon classement se lit dans le terme d'estimation : si peu de coefficients bien choisis suffisent, l'infimum sur $\bD_n$ est atteint en petite dimension, là où les règles de $\cD_n^{(d)}$ apprennent vite. La \cref{sec-proprietes} étudie la qualité du classement lui-même.
\end{enumerate}
\end{remarque}

\subsection{Le terme d'approximation}

\begin{lemme}\label{lem-borel}
Soit $\cH$ un espace de Hilbert séparable de base orthonormée $(\psi_j)_{j\ge1}$. La tribu borélienne de $\cH$ est engendrée par les formes linéaires $x\mapsto\langle x,\psi_j\rangle$, $j\ge1$.
\end{lemme}

\begin{proof}
Notons $\mathcal{G}$ la tribu engendrée par les coordonnées. Chaque coordonnée étant continue, $\mathcal{G}$ est contenue dans la tribu borélienne. Réciproquement, pour $a\in\cH$, l'application $x\mapsto\|x-a\|^2=\sum_j(\langle x,\psi_j\rangle-\langle a,\psi_j\rangle)^2$ est limite simple de fonctions $\mathcal{G}$-mesurables, donc $\mathcal{G}$-mesurable, et les boules ouvertes sont dans $\mathcal{G}$. Dans un espace métrique séparable, tout ouvert est réunion dénombrable de boules ouvertes (de centres dans une partie dénombrable dense et de rayons rationnels). Donc tout ouvert est dans $\mathcal{G}$.
\end{proof}

\begin{lemme}[{\citealp[lemme 2.1]{berlinet2008functional}}]\label{lem-approximation}
$L^*_p$ est décroissante en $J$ et $L^*_p-L^*\to0$ lorsque $J\to\infty$.
\end{lemme}

\begin{proof}
Notons $\cF_J\eqdef\sigma(\bX)=\sigma(\langle X,\psi_j\rangle,\ j\le2^J)$. C'est une filtration : les coordonnées de niveau $J$ font partie de celles de niveau $J+1$. D'après \eqref{eq-bayes-T}, $L^*_p=\E\min\{\eta_J,1-\eta_J\}$ avec $\eta_J\eqdef\E[Y\mid\cF_J]=\E[\eta(X)\mid\cF_J]$.

\emph{Décroissance.} Toute règle fonction de $\bX^{(J)}$ est une règle fonction de $\bX^{(J+1)}$, donc $L^*_{2^{J+1}}\le L^*_{2^J}$.

\emph{Convergence.} $(\eta_J)_J$ est une martingale bornée pour la filtration $(\cF_J)$. Par le théorème de convergence de Lévy, $\eta_J\to\E[\eta(X)\mid\cF_\infty]$ presque sûrement et dans $L^1$, où $\cF_\infty=\sigma\bigl(\bigcup_J\cF_J\bigr)=\sigma(\langle X,\psi_j\rangle,\ j\ge1)$. Par le \cref{lem-borel}, $\cF_\infty=\sigma(X)$, donc $\E[\eta(X)\mid\cF_\infty]=\eta(X)$. Comme $a\mapsto\min\{a,1-a\}$ est $1$-lipschitzienne,
\[
  0\le L^*_p-L^*\le\E|\eta_J-\eta(X)|\xrightarrow[J\to\infty]{}0. \qedhere
\]
\end{proof}

\begin{remarque}
\citet{berlinet2008functional} renvoient pour ce lemme à \citet[théorème 32.3]{devroye1996probabilistic}. La preuve ci-dessus est autonome et n'utilise que la complétude de la base et la séparabilité de $\Lp$ ; elle vaut pour toute base orthonormée.
\end{remarque}

\subsection{Consistance}

La consistance demande de contrôler le terme d'estimation. \citet{berlinet2008functional} le majorent par le risque d'une règle consistante de $\cD_n^{(2^J)}$. Ce passage demande une précision, car la règle est appliquée aux coordonnées \emph{réordonnées par une permutation aléatoire}.

\begin{definition}[Règle équivariante]\label{def-equivariance}
Une règle $g_n:\R^p\times(\R^p\times\{0,1\})^n\to\{0,1\}$ est \emph{équivariante par permutation} si, pour toute permutation $\sigma$, tout $\bx$ et tout échantillon,
\[
  g_n\bigl(\Pi_\sigma\bx;\,(\Pi_\sigma\bx_i,y_i)_{i\le n}\bigr)=g_n\bigl(\bx;\,(\bx_i,y_i)_{i\le n}\bigr).
\]
\end{definition}

\begin{exemple}\label{ex-equivariance}
La règle des $k$ plus proches voisins pour la distance euclidienne est équivariante, les ex æquo étant départagés selon l'indice des points d'apprentissage : $\|\Pi_\sigma\bx-\Pi_\sigma\bx_i\|=\|\bx-\bx_i\|$. L'analyse discriminante quadratique l'est aussi : les moyennes et covariances empiriques deviennent $\Pi_\sigma\hat\mu_y$ et $\Pi_\sigma\hat\Sigma_y\Pi_\sigma^\top$, et les formes quadratiques et déterminants sont invariants par conjugaison par une matrice de permutation. Les arbres à coupures parallèles aux axes le sont à la règle de départage près entre coordonnées.
\end{exemple}

\begin{remarque}[Une hypothèse implicite]\label{rem-equivariance}
Dans la preuve du corollaire 2.1 de \citet{berlinet2008functional}, la majoration $\E\{\inf_{d,g}L_n(g^{(d)})\}\le\E\{L_n(g_n^{(2^J)})\}$ utilise la consistance de $g^{(2^J)}_n$ pour la loi de $(\bX,Y)$. Or la règle est appliquée à $\Pi_{\hsigma}(\bX)$, où $\hsigma$ dépend de l'échantillon ; la loi de $(\Pi_{\hsigma}(\bX),Y)$ conditionnellement à $\Dn$ varie avec $n$. L'équivariance de $g_n$ rend l'argument exact : pour $d=p$, la règle réordonnée \emph{coïncide} avec la règle d'origine, quel que soit $\hsigma$. Les règles utilisées par les auteurs (W-NN, W-QDA) sont équivariantes, donc leur conclusion est correcte ; nous rendons simplement l'hypothèse explicite.
\end{remarque}

\begin{hypothese}\label{hyp-consistance}
Soit $\cD$ une classe de lois de $(X,Y)$. Pour tout $J\ge0$, il existe dans $\cD_n^{(2^J)}$ une règle $g_n^{(J)}$ équivariante par permutation telle que, pour toute loi de $\cD$, $\E\{L_n(g_n^{(J)})\}\to L^*_{2^J}$ lorsque $n\to\infty$, $g_n^{(J)}$ étant entraînée sur $(\bX_i,Y_i)_{i\le n}$.
\end{hypothese}

\begin{hypothese}\label{hyp-validation}
$m=m_n\to\infty$ et, pour tout $J\ge0$,
\[
  \frac{\E\{\log\Sh^{(J)}_n(2m)\}}{m}\xrightarrow[n\to\infty]{}0
  \qquad\text{et}\qquad
  \frac{\log\Card(\Gamma_n)}{m}\xrightarrow[n\to\infty]{}0 .
\]
\end{hypothese}

\begin{corollaire}[Consistance ; {\citealp[corollaire 2.1]{berlinet2008functional}}, étendu]\label{cor-consistance}
Sous les \cref{hyp-consistance,hyp-validation}, pour tout réordonnancement $(\hsigma_c)_{c\in\Gamma_n}$ fonction de $\Dn$ seulement, et pour toute loi de $\cD$,
\[
  \lim_{J\to\infty}\ \limsup_{n\to\infty}\ \E\{L_{n+m}(\hg)\}=L^* .
\]
\end{corollaire}

\begin{proof}
Fixons $J$ et $c_0\in\Gamma_n$. La règle $\bx\mapsto g_n^{(J)}\bigl(\Pi_{\hsigma_{c_0}}\bx;(\Pi_{\hsigma_{c_0}}\bX_i,Y_i)_{i\le n}\bigr)$ appartient à $\bD_n$ (terme $d=p$ de \eqref{eq-collection}). Par équivariance elle est égale à $\bx\mapsto g_n^{(J)}(\bx;(\bX_i,Y_i)_{i\le n})$. Donc
\[
  \E\Bigl\{\inf_{g\in\bD_n}L_n(g)\Bigr\}-L^*_p\le\E\{L_n(g_n^{(J)})\}-L^*_p\xrightarrow[n\to\infty]{}0
\]
par la \cref{hyp-consistance}. Pour le terme de sélection, posons $S\eqdef\Card(\Gamma_n)\,\Sh^{(J)}_n(2m)\ge1$. On a $1/\sqrt{(m/2)\log(4S)}\le1/\sqrt{(m/2)\log4}\to0$ et, par l'inégalité de Jensen appliquée à la racine carrée,
\[
  \E\sqrt{\frac{8\log(4S)}{m}}\le\sqrt{\frac{8\bigl(\log4+\log\Card(\Gamma_n)+\E\log\Sh^{(J)}_n(2m)\bigr)}{m}}\xrightarrow[n\to\infty]{}0
\]
par la \cref{hyp-validation}. Le \cref{thm-oracle} donne alors $\limsup_n\E\{L_{n+m}(\hg)\}-L^*\le L^*_p-L^*$, qui tend vers $0$ quand $J\to\infty$ par le \cref{lem-approximation}. Enfin $\E\{L_{n+m}(\hg)\}\ge L^*$ pour tout $n$.
\end{proof}

\begin{exemple}[Plus proches voisins]\label{ex-nn}
Si $\cD_n^{(d)}$ contient les règles des $k$ plus proches voisins pour $1\le k\le n$, alors $\cD_n^{(d)}$ a au plus $n$ éléments et $\Sh_{\cC_n^{(d,c)}}(2m)\le n$, d'où $\Sh^{(J)}_n(2m)\le2^Jn$. La \cref{hyp-validation} se réduit à $J/m\to0$, $\log n/m\to0$ et $\log\Card(\Gamma_n)/m\to0$. La règle $k_n$-PPV avec $k_n\to\infty$ et $k_n/n\to0$ est universellement consistante dans $\R^p$ \citep{stone1977consistent} et équivariante (\cref{ex-equivariance}). La \cref{hyp-consistance} est donc satisfaite pour \emph{toutes} les lois : la procédure est universellement consistante au sens du \cref{cor-consistance}, avec la règle de l'énergie comme avec la règle d'usage. On peut prendre $m$ égal à une petite fraction de $n$.
\end{exemple}

\begin{exemple}[Arbres]\label{ex-arbres}
Si $\cD_n^{(d)}$ contient les arbres binaires à coupures perpendiculaires aux axes, avec au plus $k$ nœuds internes, alors $\Sh_{\cC_n^{(d,c)}}(2m)\le(1+d(n+2m))^k$ \citep[chapitre 20]{devroye1996probabilistic} et $\Sh^{(J)}_n(2m)\le2^J(1+2^J(n+2m))^k$.
\end{exemple}

\subsection{Application à la règle d'usage}

\begin{theoreme}[Garanties de la règle d'usage]\label{thm-oracle-usage}
Soit $\Gamma\subset\,]0,\sqrt p[$ fini et $\hsigma_c$ la permutation de la \cref{def-usage}. La règle $\hg$ obtenue par \eqref{eq-selection} vérifie l'inégalité oracle du \cref{thm-oracle} avec ce $\Gamma$.
\end{theoreme}

\begin{proof}
$\hphi_j(c)$ est une fonction mesurable de $\bX_1,\dots,\bX_n$ (composée de la normalisation, mesurable, et d'indicatrices d'ensembles boréliens), donc $\hsigma_c$ est une fonction mesurable de $\Dn$ seulement. On applique le \cref{thm-oracle}.
\end{proof}

\begin{corollaire}\label{cor-consistance-usage}
Sous les \cref{hyp-consistance,hyp-validation}, la procédure fondée sur la règle d'usage avec seuil global est consistante au sens du \cref{cor-consistance}. Avec les plus proches voisins et une grille $\Gamma$ fixée, elle est universellement consistante dès que $J/m\to0$ et $\log n/m\to0$.
\end{corollaire}

Le point essentiel de ce corollaire est la séparation entre sélection des coefficients et classification. Si l'on classait les courbes normalisées, le terme d'approximation changerait, et la consistance pourrait être perdue.

\begin{proposition}[La normalisation peut détruire l'information]\label{prop-normalisation-perte}
Notons $\tilde L^*_p$ l'erreur de Bayes de la prédiction de $Y$ à partir de $\tX=\bX/\|\bX\|$. Alors $\tilde L^*_p\ge L^*_p$, et l'inégalité peut être stricte pour tout $J$ : il existe une loi de $(X,Y)$ telle que $L^*=L^*_p=0$ pour tout $J$ assez grand, et $\tilde L^*_p=\min\{\pi,1-\pi\}$ pour tout $J$, où $\pi=\PP(Y=1)$.
\end{proposition}

\begin{proof}
$\tX$ est une fonction mesurable de $\bX$, d'où $\tilde L^*_p\ge L^*_p$ (\cref{eq-bayes-T}). Soit $h\in\Lp$ de norme $1$ et $X=(1+Y)h$. Dès que $P_Jh\neq0$, l'application $\bX\mapsto\|\bX\|$ vaut $(1+Y)\|P_Jh\|$ et détermine $Y$, donc $L^*_p=0$. En revanche $\tX=T_J(h)/\|T_J(h)\|$ est déterministe, indépendant de $Y$, et la meilleure règle prédit la classe majoritaire.
\end{proof}

\begin{remarque}[Plus de deux classes]\label{rem-multiclasse}
Les données de parole comportent cinq classes. L'erreur de Bayes devient $\E[1-\max_k\eta_k(X)]$ avec $\eta_k(x)=\PP(Y=k\mid X=x)$, et le \cref{lem-approximation} se démontre de la même façon avec la fonction $1$-lipschitzienne (pour la norme $\ell^1$) $a\mapsto1-\max_ka_k$. Le \cref{lem-selection} et la \cref{prop-hoeffding} valent mot pour mot. Pour les familles infinies, il faut remplacer le coefficient d'éclatement par le nombre d'étiquetages distincts que la famille induit sur $N$ points.
\end{remarque}

\section{Propriétés propres à la règle d'usage}\label{sec-proprietes}

L'inégalité oracle est la même pour tous les réordonnancements. Ce qui distingue la règle d'usage, ce sont les propriétés du classement lui-même : sa stabilité, sa robustesse et son interprétation.

\subsection{Invariance et parcimonie}

\begin{proposition}\label{prop-proprietes-usage}
Soit $c>0$.
\begin{enumerate}[label=(\roman*)]
  \item \emph{Invariance d'échelle.} Si l'on remplace chaque $\bX_i$ par $a_i\bX_i$ avec $a_i\neq0$, les fréquences $\hphi_j(c)$ et la permutation $\hsigma_c$ sont inchangées. La règle de l'énergie n'a pas cette propriété.
  \item \emph{Parcimonie.} Toute courbe utilise strictement moins de $p/c^2$ fonctions de base :
  \[
    \sum_{j=1}^ps_j(\bx;c)<\frac{p}{c^2}\qquad\text{pour tout }\bx\in\R^p .
  \]
  Par conséquent $\sum_j\hphi_j(c)<p/c^2$ et, pour tout $\alpha\in\,]0,1]$,
  \[
    \Card\bigl\{j:\hphi_j(c)\ge\alpha\bigr\}<\frac{p}{\alpha c^2}.
  \]
\end{enumerate}
\end{proposition}

\begin{proof}
(i) $\widetilde{a_i\bX_i}=\operatorname{sgn}(a_i)\tX_i$ et $s_j$ ne dépend que de $|\tx_j|$. Pour l'énergie, multiplier $\bX_1$ par $a\to\infty$ fait coïncider le classement avec celui des $|X_{1j}|$.

(ii) Si $\bx=0$, aucun indice n'est utilisé. Sinon, en notant $U=\{j:s_j(\bx;c)=1\}$,
\[
  1=\sum_{j=1}^p\tx_j^2\ \ge\ \sum_{j\in U}\tx_j^2\ >\ \Card(U)\,\frac{c^2}{p}
\]
si $U\neq\emptyset$, d'où $\Card(U)<p/c^2$. En moyennant sur les courbes, $\sum_j\hphi_j(c)=\frac1n\sum_i\Card(U_i)<p/c^2$. Enfin $\alpha\,\Card\{j:\hphi_j\ge\alpha\}\le\sum_j\hphi_j$.
\end{proof}

Le point (ii) donne une lecture directe du paramètre : avec $c=2$ et $p=256$, chaque courbe utilise moins de $64$ coefficients, et moins de $128$ indices sont utilisés par au moins la moitié des courbes.

\subsection{Stabilité : concentration sans hypothèse de moment}

On définit la fréquence d'utilisation théorique
\[
  \varphi_j(c)\eqdef\PP\Bigl\{|\tX_j|>\frac{c}{\sqrt p}\Bigr\},\qquad j=1,\dots,p,
\]
et l'on note $\varphi_{(1)}(c)\ge\dots\ge\varphi_{(p)}(c)$ ses valeurs ordonnées. Lorsque l'écart
\[
  \gamma_d(c)\eqdef\varphi_{(d)}(c)-\varphi_{(d+1)}(c)
\]
est strictement positif, l'ensemble $S_d(c)$ des $d$ indices de plus grande fréquence théorique est bien défini : c'est la représentation commune \emph{théorique} en dimension $d$.

\begin{theoreme}[Stabilité de la représentation commune]\label{thm-stabilite}
Pour tout $c>0$, tout $n\ge1$ et tout $t>0$, sans aucune hypothèse sur la loi de $X$,
\[
  \PP\Bigl\{\max_{1\le j\le p}\bigl|\hphi_j(c)-\varphi_j(c)\bigr|\ge t\Bigr\}\le2p\,e^{-2nt^2}.
\]
Si de plus $\gamma_d(c)>0$, alors
\[
  \PP\bigl\{\hS_d(c)\neq S_d(c)\bigr\}\le2p\,\exp\Bigl(-\frac{n\,\gamma_d(c)^2}{2}\Bigr).
\]
En particulier, si $p=p_n$ et $\gamma_d(c)=\gamma_{d,n}(c)$ dépendent de $n$ avec $n\gamma_{d,n}(c)^2/\log p_n\to\infty$, alors $\PP\{\hS_d(c)=S_d(c)\}\to1$ ; si $\sum_np_ne^{-n\gamma_{d,n}(c)^2/2}<\infty$, alors $\hS_d(c)=S_d(c)$ pour tout $n$ assez grand, presque sûrement.
\end{theoreme}

\begin{proof}
Pour $j$ fixé, $s_j(\bX_1;c),\dots,s_j(\bX_n;c)$ sont des variables de Bernoulli i.i.d. de paramètre $\varphi_j(c)$ et $\hphi_j(c)$ est leur moyenne. L'inégalité de Hoeffding donne $\PP\{|\hphi_j-\varphi_j|\ge t\}\le2e^{-2nt^2}$ et une borne d'union sur $j$ donne la première inégalité.

Supposons $\max_j|\hphi_j-\varphi_j|<\gamma_d/2$. Pour $j\in S_d$ et $k\notin S_d$,
\[
  \hphi_j>\varphi_j-\frac{\gamma_d}{2}\ge\varphi_{(d)}-\frac{\gamma_d}{2}=\varphi_{(d+1)}+\frac{\gamma_d}{2}\ge\varphi_k+\frac{\gamma_d}{2}>\hphi_k .
\]
Chacun des $d$ indices de $S_d$ a donc une fréquence empirique strictement supérieure à celle de chaque indice hors de $S_d$ : les $d$ premiers rangs de $\hsigma_c$ sont exactement $S_d$, quelle que soit la règle de départage. Ainsi $\{\hS_d\ne S_d\}\subset\{\max_j|\hphi_j-\varphi_j|\ge\gamma_d/2\}$, et la première inégalité avec $t=\gamma_d/2$ donne la seconde. Les deux dernières affirmations en découlent directement, la dernière par le lemme de Borel--Cantelli.
\end{proof}

\begin{remarque}[Comparaison avec la règle de l'énergie]\label{rem-energie-moments}
Le pendant théorique de la règle de l'énergie est le classement des $\E X_j^2$. Il n'est défini que si ces moments sont finis. Sous cette seule hypothèse, la loi forte des grands nombres assure que l'ensemble sélectionné par l'énergie converge presque sûrement lorsque l'écart théorique correspondant est positif, mais sans aucune vitesse : une borne exponentielle analogue à celle du \cref{thm-stabilite} demande par exemple que les $X_j^2$ soient bornées ou sous-exponentielles, avec des constantes qui dépendent de la loi. Si les amplitudes ont une queue lourde, l'énergie empirique diverge et son classement n'a pas de cible. Par exemple, si $X=R\,h$ avec $h\in\Lp$ fixée, $\PP(R=0)=0$ et $\E R^2=\infty$, alors $\frac1n\sum_iX_{ij}^2\to\infty$ presque sûrement pour tout $j$ tel que $\langle h,\psi_j\rangle\neq0$. À l'inverse, $\tX=\operatorname{sgn}(R)\,T_J(h)/\|T_J(h)\|$, si bien que $\hphi_j(c)=\varphi_j(c)\in\{0,1\}$ exactement, pour tout $n$. La fréquence d'utilisation est toujours définie et la borne du \cref{thm-stabilite} est universelle : elle ne dépend que de $n$, $p$ et de l'écart $\gamma_d(c)$.
\end{remarque}

\subsection{Robustesse : influence bornée d'une courbe aberrante}

On compare maintenant deux échantillons d'apprentissage $\Dn$ et $\Dn'$ de même taille $n$ qui diffèrent par au plus $k$ courbes : $\bX'_i=\bX_i$ sauf pour au plus $k$ indices $i$, les courbes modifiées étant arbitraires. On note avec un prime les quantités calculées sur $\Dn'$, et $\hat\gamma_d(c)=\hphi_{(d)}(c)-\hphi_{(d+1)}(c)$ l'écart empirique.

\begin{proposition}[Robustesse de la règle d'usage]\label{prop-robustesse}
Pour tout $c>0$ :
\begin{enumerate}[label=(\roman*)]
  \item $\max_j|\hphi'_j(c)-\hphi_j(c)|\le k/n$ ;
  \item si $\hat\gamma_d(c)>2k/n$, alors $\hS'_d(c)=\hS_d(c)$ ;
  \item si $j_0\notin\hS_d(c)$ et $\hphi_{(d)}(c)-\hphi_{j_0}(c)>2k/n$, alors $j_0\notin\hS'_d(c)$ : pour faire entrer dans la représentation commune un indice de fréquence $\hphi_{j_0}$, il faut modifier au moins $\frac n2\bigl(\hphi_{(d)}-\hphi_{j_0}\bigr)$ courbes.
\end{enumerate}
\end{proposition}

\begin{proof}
(i) $\hphi'_j-\hphi_j=\frac1n\sum_i\bigl(s_j(\bX'_i;c)-s_j(\bX_i;c)\bigr)$ ; au plus $k$ termes sont non nuls et chacun est dans $[-1,1]$.
(ii) Même argument que dans la preuve du \cref{thm-stabilite}, avec $\hphi$ à la place de $\varphi$ et $\hphi'$ à la place de $\hphi$ : pour $j\in\hS_d$ et $l\notin\hS_d$, $\hphi'_j\ge\hphi_{(d)}-k/n>\hphi_{(d+1)}+k/n\ge\hphi_l+k/n\ge\hphi'_l$.
(iii) Pour chacun des $d$ indices $j\in\hS_d$, $\hphi'_j\ge\hphi_{(d)}-k/n>\hphi_{j_0}+k/n\ge\hphi'_{j_0}$. Au moins $d$ indices sont donc strictement devant $j_0$ pour $\hphi'$.
\end{proof}

\begin{proposition}[Point de rupture de la règle de l'énergie]\label{prop-rupture-energie}
Pour tout échantillon $\Dn$ et tout indice $j_0$, il existe un échantillon $\Dn'$ qui diffère de $\Dn$ par une seule courbe et pour lequel $j_0$ est classé premier par la règle de l'énergie. Plus généralement, en modifiant $k$ courbes, on peut placer $k$ indices arbitraires aux $k$ premiers rangs.
\end{proposition}

\begin{proof}
Soit $A^2>\sum_{i=1}^n\|\bX_i\|^2$ et $\bX'_1=A\,e_{j_0}$, les autres courbes étant inchangées. L'énergie de $j_0$ dans $\Dn'$ est au moins $A^2$, celle de tout $j\neq j_0$ est au plus $\sum_{i\ge2}X_{ij}^2\le\sum_i\|\bX_i\|^2<A^2$. Pour $k$ indices, on remplace $k$ courbes par $Ae_{j_1},\dots,Ae_{j_k}$.
\end{proof}

Les \cref{prop-robustesse,prop-rupture-energie} formalisent la motivation de la règle d'usage : dans le langage de la statistique robuste, le point de rupture de la représentation commune est de $1/n$ pour l'énergie, et de l'ordre de $\hat\gamma_d/2$ pour l'usage. Le \cref{sec-exp-contamination} vérifie ce comportement numériquement.

\subsection{Synthèse}

\begin{table}[H]
\centering
\small
\begin{tabular}{@{}>{\raggedright\arraybackslash}p{5.6cm}>{\raggedright\arraybackslash}p{3.9cm}>{\raggedright\arraybackslash}p{4.6cm}@{}}
\toprule
 & Énergie \eqref{eq-energie} & Usage à seuil global \eqref{eq-usage} \\
\midrule
Inégalité oracle et consistance & oui (\cref{thm-oracle,cor-consistance}) & oui (\cref{thm-oracle-usage,cor-consistance-usage}) \\
Coût de la sélection du seuil & aucun & $\log\Card(\Gamma)$ \\
Invariance par changement d'échelle des courbes & non & oui \\
Concentration du classement & moments d'ordre $\ge2$ requis, pas de vitesse universelle & exponentielle, sans hypothèse \\
Influence d'une courbe modifiée & non bornée & au plus $1/n$ \\
Contrôle du nombre de coefficients utilisés & non & $<p/c^2$ par courbe \\
Utilise les étiquettes & non & non \\
\bottomrule
\end{tabular}
\caption{Comparaison théorique des deux règles de représentation commune.}
\label{tab-synthese-theorie}
\end{table}
% ===== fin chapitres/05_theorie.tex =====

% ===== début chapitres/06_experiences.tex =====
\chapter{Expériences numériques}\label{chap-experiences}

% ===== début tables/classements_valeurs.tex =====
\newcommand{\JacEnergieUne}{0{,}33}
\newcommand{\JacUsageUne}{1{,}00}
\newcommand{\NbCoeffsPhoneme}{256}
% ===== fin tables/classements_valeurs.tex =====


Ce chapitre poursuit trois buts : reproduire les résultats empiriques de \citet{berlinet2008functional}, comparer la règle de l'énergie et la règle d'usage à seuil global, et vérifier sur des données les propriétés établies au \cref{chap-theorie}.

\section{Protocole}\label{sec-protocole}

Nous suivons le protocole de l'article, qui est aussi celui couvert par la théorie (\cref{sec-decoupage}). Pour chaque répétition, les données sont partagées aléatoirement, de façon stratifiée, en une séquence d'apprentissage de taille $n$, une séquence de validation de taille $m$ et une séquence de test. Le classement des coefficients n'utilise que la séquence d'apprentissage ; le triplet (seuil, dimension, règle) est choisi sur la séquence de validation par \eqref{eq-selection}, avec $d\le d_{\max}=64$ ; l'erreur est mesurée sur la séquence de test. Les deux règles sont évaluées sur les \emph{mêmes} partitions, et nous les comparons par un test des rangs signés de Wilcoxon apparié. Sauf mention contraire, nous faisons $R=100$ répétitions, comme l'article, et nous rapportons la moyenne de l'erreur de test avec son erreur standard entre parenthèses.

La sélection est faite séparément dans chaque famille (W-NN, W-QDA, W-CART), comme dans les tableaux de l'article, et conjointement sur les trois familles (ligne « Toutes »). Pour la règle d'usage, $\Gamma=\{0{,}5;\,1;\,1{,}5;\,2;\,3\}$.

\paragraph{Stabilité du classement.} Pour mesurer la variabilité de la représentation commune d'un échantillon à l'autre, nous calculons, pour chaque répétition, l'ensemble des dix premiers indices du classement (au seuil $c=1$ pour la règle d'usage), puis l'indice de Jaccard moyen entre toutes les paires de répétitions. Une valeur de $1$ signifie que toutes les partitions sélectionnent les mêmes dix coefficients.

\begin{figure}[ht]
\centering
\includegraphics[width=\textwidth]{donnees.pdf}
\caption{Les trois jeux de données. En haut : un log-périodogramme par classe de phonème ; quatre ECG de chaque classe dans le domaine temporel et dans le domaine spectral. En bas : trois courbes du modèle simulé.}
\label{fig-donnees}
\end{figure}

\section{Données de parole : reproduction de l'article}\label{sec-exp-phoneme}

Les données \emph{phoneme} \citep{hastie1995penalized} comportent $4509$ log-périodogrammes de longueur $256$, répartis en cinq classes de phonèmes (\cref{fig-donnees}). Comme dans l'article, $n=m=250$ et le test porte sur les $4009$ courbes restantes ; $J=8$. Conformément à la \cref{rem-phoneme-spectral}, les courbes sont utilisées telles quelles.

\begin{table}[ht]
\centering
\small
% ===== début tables/phoneme_bases.tex =====
\begin{tabular}{@{}llcccccc@{}}
\toprule
Famille & Règle & \texttt{db1} & \texttt{db2} & \texttt{db4} & \texttt{db6} & \texttt{db8} & \texttt{db10} \\
\midrule
\multirow{3}{*}{W-NN} & Article & 0{,}111 & 0{,}110 & 0{,}112 & 0{,}114 & 0{,}113 & 0{,}112 \\
 & Énergie & \textbf{0{,}109} & 0{,}110 & \textbf{0{,}111} & 0{,}112 & \textbf{0{,}110} & \textbf{0{,}111} \\
 & Usage & \textbf{0{,}109} & \textbf{0{,}109} & \textbf{0{,}111} & \textbf{0{,}111} & \textbf{0{,}110} & \textbf{0{,}111} \\
\addlinespace
\multirow{3}{*}{W-QDA} & Article & 0{,}097 & 0{,}102 & 0{,}108 & 0{,}108 & 0{,}113 & 0{,}115 \\
 & Énergie & 0{,}095 & 0{,}102 & 0{,}107 & \textbf{0{,}107} & 0{,}112 & \textbf{0{,}110} \\
 & Usage & \textbf{0{,}094} & \textbf{0{,}101} & \textbf{0{,}106} & \textbf{0{,}107} & \textbf{0{,}111} & \textbf{0{,}110} \\
\addlinespace
\multirow{3}{*}{W-CART} & Article & 0{,}112 & 0{,}130 & 0{,}162 & 0{,}159 & 0{,}163 & 0{,}185 \\
 & Énergie & \textbf{0{,}128} & 0{,}134 & 0{,}155 & 0{,}152 & 0{,}174 & 0{,}166 \\
 & Usage & \textbf{0{,}128} & \textbf{0{,}132} & \textbf{0{,}152} & \textbf{0{,}148} & \textbf{0{,}168} & \textbf{0{,}164} \\
\bottomrule
\end{tabular}
% ===== fin tables/phoneme_bases.tex =====

\caption{Données \emph{phoneme} : erreur de test moyenne sur $100$ partitions pour six bases de Daubechies. La ligne « Article » reprend le tableau~1 de \citet{berlinet2008functional}. En gras, la meilleure des deux règles (à $10^{-3}$ près).}
\label{tab-phoneme-bases}
\end{table}

\paragraph{Reproduction.} Le \cref{tab-phoneme-bases} reproduit le tableau~1 de l'article. Pour W-NN et W-QDA, nos erreurs avec la règle de l'énergie coïncident avec les valeurs publiées à $0{,}005$ près, pour toutes les bases, et l'on retrouve la dégradation de W-QDA lorsque le nombre de moments nuls augmente. Pour W-CART, les écarts atteignent $0{,}02$ dans les deux sens ; l'article utilise la bibliothèque R \texttt{tree}, dont l'algorithme de construction et d'élagage diffère de celui de scikit-learn.

\begin{table}[ht]
\centering
\small
% ===== début tables/phoneme_resume.tex =====
\begin{tabular}{@{}llccccc@{}}
\toprule
Famille & Règle & Erreur test & $\hat d$ médian & Stabilité top-10 & Écart & $p$-valeur \\
\midrule
\multirow{2}{*}{W-NN} & Énergie & 0{,}111 (0{,}001) & 15 & 0{,}95 & \multirow{2}{*}{$+0{,}000$} & \multirow{2}{*}{0{,}876} \\
 & Usage & 0{,}111 (0{,}001) & 15 & 0{,}90 &  &  \\
\addlinespace
\multirow{2}{*}{W-QDA} & Énergie & 0{,}107 (0{,}001) & 9 & 0{,}95 & \multirow{2}{*}{$-0{,}001$} & \multirow{2}{*}{0{,}219} \\
 & Usage & 0{,}106 (0{,}001) & 9 & 0{,}90 &  &  \\
\addlinespace
\multirow{2}{*}{W-CART} & Énergie & 0{,}155 (0{,}002) & 18 & 0{,}95 & \multirow{2}{*}{$-0{,}003$} & \multirow{2}{*}{0{,}093} \\
 & Usage & 0{,}152 (0{,}002) & 16 & 0{,}90 &  &  \\
\addlinespace
\multirow{2}{*}{Toutes} & Énergie & 0{,}107 (0{,}001) & 12 & 0{,}95 & \multirow{2}{*}{$+0{,}001$} & \multirow{2}{*}{0{,}938} \\
 & Usage & 0{,}108 (0{,}001) & 11 & 0{,}90 &  &  \\
\bottomrule
\end{tabular}
% ===== fin tables/phoneme_resume.tex =====

\caption{Données \emph{phoneme}, base \texttt{db4}. Erreur de test (erreur standard), dimension sélectionnée, stabilité du classement entre partitions, écart moyen apparié (usage moins énergie) et $p$-valeur du test de Wilcoxon apparié.}
\label{tab-phoneme-resume}
\end{table}

\paragraph{Énergie contre usage.} Sur ces données, les deux règles sont indiscernables (\cref{tab-phoneme-resume}) : les écarts appariés sont inférieurs à $0{,}004$ et aucun n'est significatif au niveau $5\,\%$. Les dimensions sélectionnées sont comparables. Ce résultat est attendu : les log-périodogrammes de parole sont des courbes régulières, d'amplitudes comparables et sans valeur aberrante marquée, et les deux critères désignent essentiellement les mêmes coefficients grossiers (\cref{fig-classements}).

La stabilité du classement est légèrement plus faible pour la règle d'usage. Ce n'est pas une contradiction avec le \cref{thm-stabilite} : la \cref{fig-classements} montre qu'une dizaine de coefficients sont utilisés par la quasi-totalité des courbes ($\hphi_j\approx1$). Les dix premiers rangs tombent donc dans un plateau où l'écart $\gamma_{10}$ est presque nul, régime dans lequel le théorème ne garantit rien et où de petites fluctuations permutent les indices.

\begin{figure}[ht]
\centering
\includegraphics[width=\textwidth]{classements.pdf}
\caption{Données \emph{phoneme} ($250$ courbes, \texttt{db4}). Scores triés des deux règles, avant (trait plein) et après (tirets) l'ajout d'un pic à une seule courbe. Le pic modifie tout le profil d'énergie ; le profil d'usage est inchangé. Indice de Jaccard entre les dix premiers indices avant et après contamination : \JacEnergieUne{} pour l'énergie, \JacUsageUne{} pour l'usage.}
\label{fig-classements}
\end{figure}

\section{Données simulées}\label{sec-exp-simulation}

Nous reprenons le modèle de la section 3.2 de l'article :
\[
  X(t)=\sin(F^1\pi t)\,f_{\mu,\sigma}(t)+\sin(F^2\pi t)\,f_{1-\mu,\sigma}(t)+\varepsilon_t ,
\]
où $f_{\mu,\sigma}$ est la densité gaussienne de moyenne $\mu$ et de variance $\sigma^2$, $F^1,F^2\sim\mathcal U[50,150]$, $\mu\sim\mathcal U[0{,}1;\,0{,}4]$, $\sigma^2\sim\mathcal U[0;\,0{,}005]$ et les $\varepsilon_t$ sont i.i.d. $\mathcal N(0,0{,}5^2)$. L'étiquette vaut $Y=0$ si $\mu\le0{,}25$ (bouffées éloignées) et $Y=1$ sinon. Les courbes sont échantillonnées en $1024$ points, $J=10$, $n=m=50$ et le test porte sur $500$ courbes indépendantes.

\begin{table}[ht]
\centering
\small
% ===== début tables/simulation_bases.tex =====
\begin{tabular}{@{}llcccccc@{}}
\toprule
Famille & Règle & \texttt{db1} & \texttt{db2} & \texttt{db4} & \texttt{db6} & \texttt{db8} & \texttt{db10} \\
\midrule
\multirow{3}{*}{W-NN} & Article & 0{,}146 & 0{,}143 & 0{,}146 & 0{,}156 & 0{,}155 & 0{,}159 \\
 & Énergie & \textbf{0{,}136} & 0{,}141 & \textbf{0{,}141} & \textbf{0{,}146} & \textbf{0{,}146} & \textbf{0{,}147} \\
 & Usage & 0{,}141 & \textbf{0{,}137} & 0{,}151 & 0{,}149 & 0{,}147 & 0{,}154 \\
\addlinespace
\multirow{3}{*}{W-QDA} & Article & 0{,}078 & 0{,}082 & 0{,}085 & 0{,}082 & 0{,}085 & 0{,}084 \\
 & Énergie & \textbf{0{,}078} & \textbf{0{,}077} & \textbf{0{,}073} & \textbf{0{,}073} & \textbf{0{,}077} & \textbf{0{,}077} \\
 & Usage & 0{,}086 & 0{,}092 & 0{,}092 & 0{,}093 & 0{,}097 & 0{,}097 \\
\addlinespace
\multirow{3}{*}{W-CART} & Article & 0{,}185 & 0{,}174 & 0{,}170 & 0{,}177 & 0{,}179 & 0{,}161 \\
 & Énergie & \textbf{0{,}132} & 0{,}136 & 0{,}132 & \textbf{0{,}133} & 0{,}132 & \textbf{0{,}140} \\
 & Usage & 0{,}137 & \textbf{0{,}131} & \textbf{0{,}126} & 0{,}137 & \textbf{0{,}131} & 0{,}141 \\
\bottomrule
\end{tabular}
% ===== fin tables/simulation_bases.tex =====

\caption{Données simulées : erreur de test moyenne sur $100$ répétitions. La ligne « Article » reprend le tableau~3 de \citet{berlinet2008functional}.}
\label{tab-simulation-bases}
\end{table}

\begin{table}[ht]
\centering
\small
% ===== début tables/simulation_resume.tex =====
\begin{tabular}{@{}llccccc@{}}
\toprule
Famille & Règle & Erreur test & $\hat d$ médian & Stabilité top-10 & Écart & $p$-valeur \\
\midrule
\multirow{2}{*}{W-NN} & Énergie & 0{,}141 (0{,}003) & 8 & 0{,}11 & \multirow{2}{*}{$+0{,}010$} & \multirow{2}{*}{0{,}054} \\
 & Usage & 0{,}151 (0{,}004) & 10 & 0{,}32 &  &  \\
\addlinespace
\multirow{2}{*}{W-QDA} & Énergie & 0{,}073 (0{,}003) & 8 & 0{,}11 & \multirow{2}{*}{$+0{,}019$} & \multirow{2}{*}{$<10^{-3}$} \\
 & Usage & 0{,}092 (0{,}004) & 9 & 0{,}32 &  &  \\
\addlinespace
\multirow{2}{*}{W-CART} & Énergie & 0{,}132 (0{,}003) & 16 & 0{,}11 & \multirow{2}{*}{$-0{,}006$} & \multirow{2}{*}{0{,}291} \\
 & Usage & 0{,}126 (0{,}003) & 15 & 0{,}32 &  &  \\
\addlinespace
\multirow{2}{*}{Toutes} & Énergie & 0{,}081 (0{,}004) & 8 & 0{,}11 & \multirow{2}{*}{$+0{,}020$} & \multirow{2}{*}{$<10^{-3}$} \\
 & Usage & 0{,}101 (0{,}004) & 11 & 0{,}32 &  &  \\
\bottomrule
\end{tabular}
% ===== fin tables/simulation_resume.tex =====

\caption{Données simulées, base \texttt{db4}. Mêmes colonnes que le \cref{tab-phoneme-resume}.}
\label{tab-simulation-resume}
\end{table}

\paragraph{Reproduction.} Avec la règle de l'énergie, nous retrouvons les conclusions de l'article (\cref{tab-simulation-bases}) : W-QDA est nettement la meilleure famille, avec une erreur d'environ $0{,}075$ contre $0{,}08$ dans l'article, et les résultats sont robustes au choix de la base. Nos erreurs pour W-CART sont plus faibles que celles de l'article, pour la même raison que précédemment.

\paragraph{Énergie contre usage.} Ici, la règle d'usage est \emph{moins bonne} en prédiction : l'écart est d'environ $+0{,}02$ pour W-QDA et pour la sélection conjointe, significatif ($p<10^{-3}$) (\cref{tab-simulation-resume}). L'explication la plus plausible tient à la structure du modèle. L'information discriminante est la \emph{position} des deux bouffées, qui varie d'une courbe à l'autre ; les coefficients fins qui la portent ne sont donc utilisés que par une petite fraction des courbes. Par construction, la règle d'usage favorise au contraire les coefficients utilisés par beaucoup de courbes, et le bruit, présent partout, en fait partie. La règle de l'énergie, elle, est attirée par les courbes à bouffées très concentrées (petit $\sigma$), dont les grands coefficients sont précisément ceux qui localisent les bouffées. Ce qui est un défaut de robustesse devient ici un avantage.

La règle d'usage produit en revanche une représentation nettement plus stable d'une répétition à l'autre (indice de Jaccard $0{,}32$ contre $0{,}11$). Enfin, le seuil sélectionné est, dans près de la moitié des répétitions, le plus grand de la grille, $c=3$ : la grille est trop étroite pour ce modèle. Le \cref{sec-exp-grille} montre qu'avec une grille plus large l'écart avec la règle de l'énergie cesse d'être significatif.

Cet exemple illustre la contrepartie annoncée au \cref{chap-methode} : la règle d'usage cherche une représentation \emph{commune}, qui décrit ce que partagent les courbes, et non la représentation la plus discriminante. Lorsque la classe dépend d'un détail rare et mobile, les deux objectifs divergent.

\section{ECG200 : domaine temporel et domaine spectral}\label{sec-exp-ecg}

Le jeu ECG200 \citep{olszewski2001generalized,dau2019ucr} contient $200$ électrocardiogrammes de longueur $96$, répartis en deux classes (battement normal ou infarctus du myocarde). Contrairement aux données de parole, ce sont de vraies séries temporelles : on peut comparer le domaine temporel au domaine spectral obtenu par le log-périodogramme (\cref{sec-periodogramme}), qui comporte $48$ ordonnées. Nous regroupons les deux échantillons d'origine et, à chaque répétition, nous tirons $n=m=50$ courbes, le test portant sur les $100$ restantes.

\begin{table}[ht]
\centering
\small
% ===== début tables/ecg200.tex =====
\begin{tabular}{@{}lllccc@{}}
\toprule
Domaine & Famille & Règle & Erreur test & $\hat d$ médian & $p$-valeur \\
\midrule
Temporel & W-NN & Énergie & 0{,}164 (0{,}005) & 13 & 0{,}063 \\
 &  & Usage & 0{,}173 (0{,}005) & 11 &  \\
 & W-QDA & Énergie & 0{,}176 (0{,}005) & 10 & 0{,}796 \\
 &  & Usage & 0{,}173 (0{,}004) & 9 &  \\
 & W-CART & Énergie & 0{,}212 (0{,}004) & 16 & 0{,}254 \\
 &  & Usage & 0{,}208 (0{,}004) & 13 &  \\
 & Toutes & Énergie & 0{,}170 (0{,}004) & 11 & 0{,}282 \\
 &  & Usage & 0{,}170 (0{,}005) & 11 &  \\
\midrule
Périodogramme & W-NN & Énergie & 0{,}224 (0{,}004) & 22 & 0{,}597 \\
 &  & Usage & 0{,}219 (0{,}004) & 17 &  \\
 & W-QDA & Énergie & 0{,}221 (0{,}004) & 7 & $<10^{-3}$ \\
 &  & Usage & 0{,}204 (0{,}004) & 6 &  \\
 & W-CART & Énergie & 0{,}206 (0{,}005) & 16 & 0{,}587 \\
 &  & Usage & 0{,}206 (0{,}005) & 12 &  \\
 & Toutes & Énergie & 0{,}210 (0{,}004) & 14 & 0{,}130 \\
 &  & Usage & 0{,}205 (0{,}004) & 11 &  \\
\bottomrule
\end{tabular}
% ===== fin tables/ecg200.tex =====

\caption{ECG200, base \texttt{db4} : erreur de test (erreur standard) sur $100$ répétitions, dimension sélectionnée et $p$-valeur du test de Wilcoxon apparié entre les deux règles.}
\label{tab-ecg}
\end{table}

\paragraph{Domaine.} Le domaine temporel est nettement meilleur : pour la sélection conjointe, l'erreur passe d'environ $0{,}17$ dans le domaine temporel à $0{,}21$ dans le domaine spectral, quelle que soit la règle (\cref{tab-ecg}). C'est la conséquence attendue de l'invariance du périodogramme par translation (\cref{sec-periodogramme}) : ce qui distingue un battement pathologique est la forme et la position de certaines ondes, information de phase que le périodogramme élimine. Pour les phonèmes, au contraire, l'information utile est la répartition de l'énergie selon les fréquences, et le domaine spectral est le bon. Le choix du domaine relève donc de la connaissance du signal, et non d'une règle générale.

\paragraph{Règles.} Dans le domaine temporel, les deux règles sont équivalentes (aucune différence significative). Dans le domaine spectral, la règle d'usage est significativement meilleure avec W-QDA (erreur $0{,}204$ contre $0{,}221$, $p<10^{-3}$) et sélectionne des dimensions plus faibles. Un périodogramme est une estimation bruitée (\cref{prop-periodogramme}) dont quelques ordonnées peuvent être très grandes ; ce contexte est favorable à une règle robuste.

\section{Robustesse à des courbes aberrantes}\label{sec-exp-contamination}

Pour vérifier les \cref{prop-robustesse,prop-rupture-energie}, nous reprenons le modèle simulé et nous ajoutons un pic étroit (gaussienne d'écart type $2$ pas, de hauteur $50$ fois le maximum absolu de l'échantillon, à une position uniforme) à $k\in\{0,1,3,5\}$ courbes de la séquence d'apprentissage, sur $n=50$. Les séquences de validation et de test ne sont pas modifiées, et les étiquettes non plus : ce sont des artefacts d'enregistrement. Pour chaque répétition, nous comparons les dix premiers indices du classement avec ceux obtenus sur l'échantillon propre, sur les mêmes données.

\begin{table}[ht]
\centering
\small
% ===== début tables/contamination.tex =====
\begin{tabular}{@{}lcccc@{}}
\toprule
Courbes contaminées (sur $n=50$) & 0 & 1 & 3 & 5 \\
\midrule
Jaccard top-10 avec l'échantillon propre -- énergie & 1{,}00 & 0{,}02 & 0{,}03 & 0{,}03 \\
Jaccard top-10 avec l'échantillon propre -- usage & 1{,}00 & 0{,}89 & 0{,}83 & 0{,}77 \\
\addlinespace
Erreur test -- énergie & 0{,}081 & 0{,}130 & 0{,}130 & 0{,}123 \\
Erreur test -- usage & 0{,}101 & 0{,}099 & 0{,}104 & 0{,}104 \\
\addlinespace
$\hat d$ médian -- énergie & 8 & 34 & 13 & 12 \\
$\hat d$ médian -- usage & 11 & 11 & 11 & 12 \\
\bottomrule
\end{tabular}
% ===== fin tables/contamination.tex =====

\caption{Contamination de la séquence d'apprentissage ($100$ répétitions, \texttt{db4}, sélection conjointe). Jaccard : recouvrement moyen des dix premiers indices avec ceux de l'échantillon propre.}
\label{tab-contamination}
\end{table}

\begin{figure}[ht]
\centering
\includegraphics[width=\textwidth]{contamination.pdf}
\caption{Effet de la contamination. À gauche : recouvrement des dix premiers indices avec ceux de l'échantillon propre. À droite : erreur de test (barres : $\pm$ une erreur standard).}
\label{fig-contamination}
\end{figure}

Le résultat est sans ambiguïté (\cref{tab-contamination,fig-contamination}). \emph{Une seule} courbe contaminée sur $50$ suffit à renouveler presque entièrement les dix premiers indices de la règle de l'énergie (Jaccard $0{,}02$) : les coefficients fins qui portent le pic passent en tête, exactement comme dans la preuve de la \cref{prop-rupture-energie}. L'erreur de test passe de $0{,}081$ à $0{,}130$ (test de Wilcoxon apparié, $p<10^{-10}$) et la dimension sélectionnée médiane de $8$ à $34$ : il faut aller chercher les coefficients utiles plus loin dans le classement. La règle d'usage n'est presque pas affectée : le recouvrement reste de $0{,}89$ avec une courbe contaminée et de $0{,}77$ avec cinq, et son erreur ne change pas de façon significative ($0{,}101$ sans contamination, $0{,}104$ avec cinq courbes, $p=0{,}24$). Dès qu'une seule courbe est contaminée, la règle d'usage devient significativement meilleure que la règle de l'énergie (écart de $-0{,}031$ pour une courbe, $p<10^{-6}$).

La \cref{prop-robustesse} explique ce comportement : chaque courbe contaminée déplace chaque fréquence d'utilisation d'au plus $1/n=0{,}02$, et le pic, qui concentre l'énergie de la courbe contaminée sur quelques coefficients, ne peut « voter » que pour un petit nombre d'indices (\cref{prop-proprietes-usage}\,(ii)).

\section{Comportement quand la taille d'échantillon augmente}\label{sec-exp-taille}

Le \cref{cor-consistance-usage} affirme que la règle d'usage conduit, comme la règle de l'énergie, à une procédure consistante. Nous l'illustrons sur le modèle simulé, avec $500$ courbes de test, $100$ répétitions et des séquences d'apprentissage et de validation de tailles $n=m\in\{25,50,100,200,400\}$.

\begin{figure}[ht]
\centering
\includegraphics[width=\textwidth]{taille_echantillon.pdf}
\caption{Modèle simulé : erreur de test (à gauche, $\pm$ une erreur standard) et dimension sélectionnée médiane (à droite) en fonction de la taille des séquences d'apprentissage et de validation.}
\label{fig-taille}
\end{figure}

Les deux erreurs décroissent vers $0$ et l'écart entre les règles se referme (\cref{fig-taille}) : $0{,}159$ contre $0{,}143$ pour $n=25$, puis $0{,}030$ contre $0{,}028$ pour $n=200$ et $0{,}018$ pour les deux règles pour $n=400$. La dimension sélectionnée croît avec $n$, conformément à l'arbitrage décrit par le \cref{thm-oracle} : avec plus de données, les règles de $\cD_n^{(d)}$ apprennent bien en dimension plus grande, et la sélection va chercher les coefficients plus fins. L'avantage de la règle de l'énergie observé au \cref{sec-exp-simulation} (avec la grille par défaut) est donc un effet de petit échantillon : il porte sur la vitesse, pas sur la limite.

\begin{figure}[ht]
\centering
\includegraphics[width=\textwidth]{dimensions.pdf}
\caption{Distribution de la dimension sélectionnée $\hat d$ sur $100$ répétitions (\texttt{db4}), pour chaque famille et chaque règle.}
\label{fig-dimensions}
\end{figure}

\section{Sensibilité à la grille de seuils}\label{sec-exp-grille}

Au \cref{sec-exp-simulation}, le seuil retenu était souvent le plus grand de la grille $\Gamma=\{0{,}5;\,1;\,1{,}5;\,2;\,3\}$. Nous refaisons donc les expériences, sur les mêmes partitions, avec la grille élargie $\Gamma'=\{0{,}5;\,1;\,1{,}5;\,2;\,3;\,4;\,6;\,8\}$, qui reste dans l'intervalle $]0,\sqrt p[$ pour les trois jeux de données ($\sqrt p\ge\sqrt{96}\approx9{,}8$). D'après le \cref{thm-oracle}, cela ne coûte qu'un terme en $\log\Card(\Gamma')$.

\begin{table}[ht]
\centering
\small
% ===== début tables/grille.tex =====
\begin{tabular}{@{}lccccc@{}}
\toprule
Données & Grille & Erreur test & $\hat d$ médian & $\hat c$ le plus fréquent & $p$-valeur \\
\midrule
simulation & $\{0{,}5;\dots;3\}$ & 0{,}101 (0{,}004) & 11 & 3{,}0 & $<10^{-3}$ \\
 & $\{0{,}5;\dots;8\}$ & 0{,}085 (0{,}004) & 8 & 8{,}0 &  \\
\addlinespace
phoneme & $\{0{,}5;\dots;3\}$ & 0{,}108 (0{,}001) & 11 & 0{,}5 & 1{,}000 \\
 & $\{0{,}5;\dots;8\}$ & 0{,}108 (0{,}001) & 11 & 0{,}5 &  \\
\addlinespace
ECG200 (temps) & $\{0{,}5;\dots;3\}$ & 0{,}170 (0{,}005) & 11 & 0{,}5 & 0{,}400 \\
 & $\{0{,}5;\dots;8\}$ & 0{,}173 (0{,}005) & 11 & 0{,}5 &  \\
\bottomrule
\end{tabular}
% ===== fin tables/grille.tex =====

\caption{Règle d'usage avec la grille par défaut et avec la grille élargie (\texttt{db4}, sélection conjointe, $100$ répétitions appariées). La $p$-valeur compare les deux grilles.}
\label{tab-grille}
\end{table}

Sur les données simulées, la grille élargie fait baisser l'erreur de $0{,}101$ à $0{,}085$ ($p<10^{-3}$), avec des dimensions plus faibles ; les seuils $c=6$ et $c=8$ sont retenus dans plus de la moitié des répétitions (\cref{tab-grille}). L'écart avec la règle de l'énergie ($0{,}081$) n'est alors plus significatif (écart apparié $+0{,}004$, $p=0{,}44$). Sur \emph{phoneme} et sur ECG200, rien ne change de façon significative : les grands seuils ne sont presque jamais retenus.

L'interprétation est directe. Avec $c$ grand, une courbe n'« utilise » un coefficient que s'il porte une part importante de son énergie, soit, pour $p=1024$ et $c=8$, plus de $6\,\%$ de l'énergie totale. On ne retient ainsi que les coefficients dominants de chaque courbe, qui sont ici ceux qui localisent les bouffées. La règle garde ses propriétés de robustesse, puisque chaque courbe vote toujours pour moins de $p/c^2$ indices et que son influence reste bornée par $1/n$. Nous recommandons donc de prendre une grille qui couvre tout l'intervalle $]0,\sqrt p[$, par exemple jusqu'à $c\approx\sqrt p/2$ ; c'est désormais la grille par défaut du code.
% ===== fin chapitres/06_experiences.tex =====

% ===== début chapitres/07_conclusion.tex =====
\chapter{Conclusion et perspectives}\label{chap-conclusion}

\section{Bilan}

Ce mémoire est parti de la méthode de classification de courbes par ondelettes de \citet{berlinet2008functional} et poursuivait trois objectifs.

\paragraph{Reproduction.} Nous avons redémontré en détail l'inégalité oracle, le lemme d'approximation et le corollaire de consistance de l'article, en établissant les inégalités de sélection sur lesquelles ils reposent. La relecture de la preuve a fait apparaître une hypothèse implicite, l'équivariance par permutation des règles de classification, que nous avons rendue explicite (\cref{rem-equivariance}). Sur le plan empirique, nous retrouvons les erreurs publiées sur les données de parole à quelques millièmes près pour W-NN et W-QDA, ainsi que les conclusions de l'article sur les données simulées.

\paragraph{Domaine spectral.} Nous avons formalisé l'étape de périodogramme d'après \citet{shumway2011time} : estimateur non consistant de la densité spectrale, stabilisation de la variance par le logarithme, invariance par translation. Cette dernière propriété explique nos observations : le domaine spectral est le bon pour la parole, dont l'information est fréquentielle, mais pas pour les électrocardiogrammes, dont l'information est aussi dans la position des ondes.

\paragraph{Nouvelle règle de représentation commune.} Nous avons proposé de classer les fonctions d'ondelettes selon leur fréquence d'utilisation, définie par un seuil global $\tau=c/\sqrt p$ appliqué aux courbes ramenées à la même énergie. Sur le plan théorique :
\begin{itemize}
  \item la règle hérite de l'inégalité oracle et de la consistance de l'article, le choix du seuil dans une grille finie ne coûtant qu'un terme logarithmique (\cref{thm-oracle-usage,cor-consistance-usage}) ;
  \item la normalisation doit servir à \emph{choisir} les coefficients et non à classer, sous peine de perdre l'information portée par l'amplitude (\cref{prop-normalisation-perte}) ;
  \item la représentation commune empirique converge exponentiellement vite vers sa version théorique, sans aucune hypothèse de moment (\cref{thm-stabilite}) ;
  \item une courbe modifiée déplace chaque fréquence d'au plus $1/n$, alors qu'une seule courbe suffit à imposer n'importe quel indice à la règle de l'énergie (\cref{prop-robustesse,prop-rupture-energie}).
\end{itemize}
Sur le plan empirique, la règle d'usage fait jeu égal avec la règle de l'énergie sur les données de parole et sur les ECG dans le domaine temporel. Elle est meilleure avec W-QDA sur les périodogrammes d'ECG et, surtout, beaucoup plus robuste : une seule courbe aberrante sur cinquante dégrade fortement la règle de l'énergie et laisse la règle d'usage intacte. Sur les données simulées de l'article, où l'information discriminante est portée par des coefficients rares, elle est moins performante avec une grille de seuils trop étroite ; avec une grille couvrant les grands seuils, ou lorsque la taille d'échantillon augmente, l'écart avec la règle de l'énergie n'est plus significatif.

La règle d'usage n'est donc pas une règle « plus performante » : elle définit une représentation commune plus robuste et plus facile à interpréter, dont les performances sur données propres sont comparables à celles de la règle de l'énergie à condition de bien choisir la grille de seuils. C'est le compromis visé.

\section{Limites}

\begin{itemize}
  \item Les garanties portent sur le découpage apprentissage/validation. La validation croisée, plus économe en données, n'est pas couverte.
  \item Le \cref{thm-stabilite} suppose un écart $\gamma_d(c)$ strictement positif. Sur les données de parole, plusieurs coefficients sont utilisés par presque toutes les courbes, et le classement à l'intérieur de ce plateau n'est pas stable.
  \item La théorie est écrite pour deux classes ; l'extension à plusieurs classes est esquissée (\cref{rem-multiclasse}) mais pas détaillée pour les familles infinies.
  \item Les expériences de robustesse utilisent une contamination artificielle. Il faudrait les confirmer sur des données réelles contenant des artefacts.
\end{itemize}

\section{Perspectives}

\begin{itemize}
  \item \textbf{Seuils par niveau de résolution.} Le bruit se répartit uniformément sur les coefficients d'une transformée orthogonale, mais le signal non. Un seuil $c_j/\sqrt p$ dépendant du niveau $j$ pourrait mieux séparer les deux, au prix d'une grille plus grande, dont le coût reste logarithmique dans l'inégalité oracle.
  \item \textbf{Vitesses de convergence.} Sous des hypothèses de régularité de la fonction de régression dans une boule de Besov, on pourrait chercher à quantifier le terme d'approximation $L^*_p-L^*$ en fonction de $J$ et en déduire un choix de $J$ en fonction de $n$.
  \item \textbf{Règles hybrides.} Combiner fréquence d'utilisation et énergie, par exemple en classant par énergie moyenne \emph{parmi les courbes qui utilisent le coefficient}, pourrait conserver la robustesse tout en récupérant les coefficients rares mais très informatifs.
  \item \textbf{Variante supervisée.} La variante qui classe les coefficients par l'écart de fréquence d'utilisation entre classes a été mise de côté, car elle ne définit plus une représentation commune. Son étude théorique demanderait de contrôler le sur-apprentissage du critère de sélection lui-même.
\end{itemize}
% ===== fin chapitres/07_conclusion.tex =====


\appendix
% ===== début chapitres/A_implementation.tex =====
\chapter{Implémentation et reproductibilité}\label{chap-implementation}

\section{Organisation du code}

Le code est écrit en Python et s'appuie sur PyWavelets \citep{lee2019pywavelets} pour les transformées en ondelettes et sur scikit-learn \citep{pedregosa2011scikit} pour les classifieurs. Il est organisé en modules :

\begin{center}
\small
\begin{tabular}{@{}p{4.6cm}p{9.6cm}@{}}
\toprule
Module & Rôle \\
\midrule
\texttt{coeffient\_compute.py} & Transformée en ondelettes périodisée orthogonale (\texttt{coeff\_matrix}, \texttt{dyadic\_level}) et log-périodogramme (\texttt{to\_periodogram}). \\
\texttt{representation.py} & Règles de représentation commune : énergie, fréquence d'utilisation, normalisation des courbes (\texttt{normalize\_rows}), seuil global (\texttt{significance\_mask}, règle \texttt{"global"}). \\
\texttt{benchmark.py} & Protocole de découpage apprentissage/validation/test répété, pour tous les réglages du \cref{chap-experiences}. \\
\texttt{simulation.py} & Modèle simulé de \citet[section 3.2]{berlinet2008functional} et contamination par des pics. \\
\texttt{report.py} & Production des tableaux et figures du mémoire à partir des résultats. \\
\texttt{model.py} & Pipeline interactif avec sélection par validation croisée à $K$ blocs, piloté par \texttt{configs/config.yaml}. \\
\bottomrule
\end{tabular}
\end{center}

\section{Choix d'implémentation}

\begin{itemize}
  \item \textbf{Transformée orthogonale.} \texttt{pywt.wavedec} est appelée avec \texttt{mode="periodization"} et un nombre de niveaux égal à la plus grande puissance de $2$ divisant la longueur du signal (\cref{prop-dwt-orthogonale,rem-mode-symmetric}). Un test vérifie la propriété de Parseval pour \texttt{db1}, \texttt{db4} et \texttt{db10} et des longueurs $96$, $256$ et $1024$.
  \item \textbf{Normalisation.} Elle n'est appliquée qu'au calcul des fréquences d'utilisation ; les classifieurs reçoivent les coefficients d'origine (\cref{prop-normalisation-perte}).
  \item \textbf{Départage des ex æquo.} Tri stable par score décroissant : à score égal, l'indice le plus grossier passe en premier.
  \item \textbf{Classifieurs.} W-NN : $k\in\{1,3,5,7,9,15,25\}$ ; W-QDA : analyse discriminante quadratique sans régularisation, pour $d$ strictement inférieur à l'effectif de la plus petite classe d'apprentissage, comme dans l'article ; W-CART : arbre de décision de scikit-learn avec ses paramètres par défaut.
  \item \textbf{Sélection.} Erreur de validation minimale, puis plus petite dimension en cas d'égalité.
  \item \textbf{Aléa.} La répétition $r$ utilise la graine $1000+r$ pour le découpage, la simulation et la contamination. Les deux règles sont évaluées sur les mêmes partitions, ce qui permet des comparaisons appariées.
\end{itemize}

\section{Tests}

Le fichier \texttt{tests/test\_representation.py} contient 22 tests, qui passent tous. Ils vérifient notamment l'orthogonalité de la transformée, la définition du seuil global, l'invariance par changement d'échelle (\cref{prop-proprietes-usage}), la borne d'influence $1/n$ (\cref{prop-robustesse}), et l'exécution de bout en bout du pipeline.

\section{Reproduire les résultats}

Depuis la racine du dépôt :
\begin{verbatim}
mamba env create -f environment.yml
conda activate functional_supervised_classification_environment
pip install -e .
python tests/test_representation.py
python -m functional_supervised_classification.benchmark all --reps 100
python -m functional_supervised_classification.report
\end{verbatim}
La deuxième commande écrit les résultats bruts dans \texttt{results/} ; la troisième produit les fichiers de \texttt{memoire/tables/} et \texttt{memoire/figures/} inclus dans ce document.
% ===== fin chapitres/A_implementation.tex =====


\cleardoublepage
\phantomsection
\addcontentsline{toc}{chapter}{Bibliographie}
\printbibliography[title={Bibliographie}]

\end{document}
